\chapter{Verilator}
\label{ch:verilator}

\chapabstract{Verilator is not a simulator in the sense you learned with
\tool{QuestaSim}: it is a compiler that turns synthesisable \SV{} into C++,
which you then build into your own program. This chapter explains what that
changes, what it buys you and what it costs you; it takes you from a
forty-line \code{main()} to a complete UVM-shaped verification environment
written in C++, with coverage, assertions and lint; and it says honestly which
parts of \SV{} the tool does not implement. Every command, every listing and
every block of output in this chapter was executed on Verilator 5.020.}

\section{What Verilator Is}
\label{sec:verilator:what}

\index{Verilator}Every simulator you have used so far reads your \VHDL{},
elaborates it into an internal data structure, and then runs an
\emph{event-driven kernel} over that structure: a queue of pending events, a
notion of delta cycles, a resolution function on every signal. \tool{ModelSim}
and \tool{QuestaSim} both work this way, whichever language you feed them.
You interact with the kernel through a prompt (\code{vsim}), a waveform window
and a Tcl interpreter.

Verilator does none of this. It reads synthesisable \SV{} and \emph{writes C++
source code} that computes the same function. That C++ is then compiled by
\tool{g++} or \tool{clang++} and linked into a program that you wrote. The
design becomes a C++ class; its ports become public data members; advancing
the simulation means calling a method.

\index{Verilator!name}The name is a contraction of \emph{Verilog} and
\emph{translator}, which is exactly what the tool does. It was written by
Wilson Snyder and has been developed in the open since 1994; it is dual
licensed under the \emph{LGPL version 3} and the \emph{Artistic License 2.0},
so you may use it on proprietary designs and link the generated code into
proprietary programs without any licence obligation on your \RTL{}. There is
no licence server, no seat count and no annual renewal. For a student that is
not a small detail: it is the difference between verifying at home and
verifying only in the laboratory.

\begin{notebox}{``Verilator is a compiler'' is the whole chapter}
If you remember one sentence, remember this one. Almost every surprise in
this chapter follows from it. There is no interactive prompt, because there is
no kernel to talk to. There is no \code{\$time} unless your program maintains
one. There is no ``restart'' command, because restarting means running your
executable again. And a syntax error is a \emph{compile} error, reported
before anything runs at all.
\end{notebox}

The flow has three stages, shown in \cref{fig:verilator:flow}. First Verilator
itself (\code{verilator}) reads the \SV{} and emits C++ into a directory
called \file{obj\_dir} by default. Then \tool{make} compiles that C++ together
with your own \file{main.cpp} and with the small Verilator runtime library
into an ordinary executable. Finally you run the executable, which may write a
waveform file in VCD or FST format.

\begin{figure}[htbp]
  \centering
  \begin{tikzpicture}[
      node distance=6mm,
      box/.style   ={draw=codeframe,fill=codebg,rounded corners=2pt,
                     minimum height=8mm,inner xsep=5pt,font=\small\sffamily},
      tool/.style  ={draw=codekey,fill=notebg,rounded corners=2pt,
                     minimum height=8mm,inner xsep=5pt,
                     font=\small\sffamily\bfseries,text=codekey},
      outb/.style  ={draw=accent,fill=tipbg,rounded corners=2pt,
                     minimum height=8mm,inner xsep=5pt,font=\small\sffamily},
      ar/.style    ={-{Latex[length=2mm]},draw=codenumber,thick}]

    \node[box]                        (sv)   {\file{cordic\_rot.sv}};
    \node[box,below=2mm of sv]        (pkg)  {\file{cordic\_pkg.sv}};
    \node[tool,right=10mm of sv]      (vl)   {\texttt{verilator -{}-cc}};
    \node[box,right=10mm of vl]       (cpp)  {\file{obj\_dir/Vcordic\_rot.cpp}};
    \node[box,below=2mm of cpp]       (hdr)  {\file{obj\_dir/Vcordic\_rot.h}};

    \node[box,below=16mm of pkg]      (main) {\file{main.cpp} (yours)};
    \node[tool,right=10mm of main]    (mk)   {\texttt{make}};
    \node[outb,right=10mm of mk]       (exe)  {\file{obj\_dir/Vcordic\_rot}};
    \node[outb,below=8mm of exe]       (vcd)  {\file{dump.vcd} / \file{.fst}};

    \draw[ar] (sv.east)   -- (vl.west);
    \draw[ar] (pkg.east)  -- (vl.west);
    \draw[ar] (vl.east)   -- (cpp.west);
    \draw[ar] (vl.east)   -- (hdr.west);
    \draw[ar] (cpp.south) |- ($(mk.north)+(0,4mm)$) -- (mk.north);
    \draw[ar] (main.east) -- (mk.west);
    \draw[ar] (mk.east)   -- (exe.west);
    \draw[ar] (exe.south) -- (vcd.north) node[midway,right,font=\scriptsize\sffamily]{run};

    \node[font=\scriptsize\itshape,below=1mm of mk] {g++, plus the Verilated runtime};
  \end{tikzpicture}
  \caption{The Verilator flow. The middle stage is an ordinary C++ build; the
  last stage is an ordinary program.}
  \label{fig:verilator:flow}
\end{figure}

\section{Verilator and QuestaSim Are Different Machines}
\label{sec:verilator:vs-questa}

This section is the one to read slowly. Everything that goes wrong later goes
wrong because a habit formed against an event-driven four-state simulator was
carried unchanged into a compiled two-state one.

\subsection{Event-driven versus statically scheduled}
\label{sec:verilator:sched}

\index{event queue}\index{scheduling!static}A \VHDL{} simulator holds, at
every point in time, a queue of pending signal updates and a set of processes
waiting on events. When a signal changes, the kernel looks up its sensitivity
list, wakes the processes that care, runs them, collects the new transactions,
and repeats until no more delta cycles are pending. That machinery is what
lets you write \code{wait for 3 ns}, model a transport delay, resolve two
drivers on a bus, and stop the simulation at an arbitrary instant to inspect
it. It also costs: every signal is an object, every update is an event,
and every event goes through the queue.

Verilator throws the queue away. At compile time it analyses the whole design
and produces a \emph{static schedule}: a straight-line C++ function that
evaluates every piece of combinational logic exactly once, in an order that
respects the data dependencies, and then commits the sequential updates. The
transformations that make this possible are the same ones a synthesis tool
performs:

\begin{itemize}
  \item \textbf{Hierarchy flattening.} Module instances are inlined. In the
        generated code the CORDIC's internal signals live in one flat
        structure named \code{Vcordic\_rot\_\_\_024root}; the module boundary
        has ceased to exist.
  \item \textbf{Constant propagation.} \code{ATAN\_ROM} is computed once by
        Verilator, at elaboration, and appears in the C++ as literals.
  \item \textbf{Dead-logic elimination.} Anything that cannot affect an output
        or a traced signal is deleted.
  \item \textbf{Bit blasting and re-packing.} Narrow signals are packed into
        machine words; a 20-bit \code{idata\_t} becomes a \code{uint32\_t}.
  \item \textbf{Activity regions.} The generated \code{eval()} is split into
        blocks guarded by ``did anything in this cone change?'' flags, so a
        quiet part of the design costs a comparison and nothing more.
\end{itemize}

The consequence for you is simple and unforgiving: \emph{the intermediate
states that an event-driven simulator would show you may not exist at all}. If
two \code{always\_comb} blocks would, in Questa, briefly disagree during a
delta cycle, Verilator has already reordered them so that the disagreement
never occurs. That is usually what you want, because a synthesised netlist
would not show the glitch either, but it means Verilator is a poor tool for
studying race conditions.

\begin{gotcha}{Verilator will not reproduce a delta-cycle race}
If your \SV{} has a genuine race -- a testbench that drives a signal with a
blocking assignment in the same time step in which the \DUT{} samples it --
Questa may show it and Verilator may not, or the other way round. Verilator
detects some of these statically and reports \code{UNOPTFLAT} or
\code{MULTIDRIVEN}; it cannot detect all of them. Race hunting belongs in an
event-driven simulator. Use clocking blocks (\refverificationbook) so that the race
cannot exist in the first place.
\end{gotcha}

\subsection{Two states, not four}
\label{sec:verilator:twostate}

\index{four-state logic}\index{X propagation}This is the single most important
warning in the chapter, so it gets a demonstration rather than an assertion.

\VHDL{}'s \code{std\_logic} has nine values; \SV{}'s \code{logic} has four:
\code{0}, \code{1}, \code{X} and \code{Z}. Verilator's internal model has
\emph{two}. Every signal is a plain C integer; there is no room for an unknown.
Consider a design with a classic reset bug -- an accumulator that is used but
never reset:

\begin{svcode}[caption={A register that is used but never reset},label={lst:verilator:resetbug}]
module resetbug (
  input  logic       clk_i,
  input  logic       rst_ni,
  input  logic [7:0] d_i,
  output logic [7:0] sum_o
);
  logic [7:0] acc;
  logic       en;

  always_ff @(posedge clk_i or negedge rst_ni) begin
    if (!rst_ni) en <= 1'b0;      // en is reset ...
    else         en <= 1'b1;
  end

  always_ff @(posedge clk_i) begin
    if (en) acc <= acc + d_i;     // ... acc is not
  end

  assign sum_o = acc;
endmodule
\end{svcode}

A four-state simulator propagates the unknown initial value of \code{acc}
through the adder and out of the module, and the testbench sees \code{X}. Here
is the same testbench run under \tool{Icarus Verilog} 12.0, which is
event-driven and four-state, and then under Verilator 5.020:

\begin{txtcode}
$ iverilog -g2012 -s tb_resetbug -o rb.vvp resetbug.sv && vvp rb.vvp
sum_o after reset + 4 clocks = xxxxxxxx (x)
VERDICT: X detected -- the reset bug is visible

$ verilator --binary --timing --top-module tb_resetbug resetbug.sv -o rb
$ ./obj_dir/rb
sum_o after reset + 4 clocks = 00000011 (3)
VERDICT: sum_o is a clean number -- bug invisible
\end{txtcode}

Verilator gives you a clean, plausible, entirely wrong number. Nothing in the
run suggests that anything is missing. A scoreboard comparing against a golden
model that also starts from zero would happily report \emph{pass}.

\begin{gotcha}{A reset bug that Questa catches can hide in Verilator}
Verilator is \emph{X-optimistic}: an unknown becomes a definite value, usually
zero, and the missing reset disappears. Four-state simulators are
\emph{X-pessimistic} in the opposite direction: an \code{X} on the select of a
multiplexer poisons the output even when both inputs are equal, which produces
false failures. Neither model is the silicon. The practical rule is: run reset
and initialisation checks in a four-state simulator, and run bulk regression in
Verilator.
\end{gotcha}

Verilator offers three switches that soften the problem. They are worth
knowing precisely, because their names are similar and their effects are not.

\begin{itemize}
  \item \code{-{}-x-assign 0|1|unique|fast} controls what an \emph{explicit}
        \code{'x} in your source becomes. The default is \code{fast}, which
        lets the optimiser choose whatever is cheapest.
  \item \code{-{}-x-initial 0|unique|fast} controls the initial value of a
        variable that has no reset. The default is \code{unique}.
  \item At \emph{run} time, \code{+verilator+rand+reset+0|1|2} overrides the
        \code{unique} choice with all-zeros, all-ones or random, and
        \code{+verilator+seed+N} picks the random seed.
\end{itemize}

Measured on the counter of \cref{lst:verilator:resetbug}:

\begin{txtcode}
$ ./obj_dir/rb                                             # default
sum_o after reset + 4 clocks = 00000011 (3)
$ ./obj_dir/rb +verilator+rand+reset+2 +verilator+seed+1
sum_o after reset + 4 clocks = 00000011 (3)
$ ./obj_dir/rb +verilator+rand+reset+2 +verilator+seed+2
sum_o after reset + 4 clocks = 00000010 (2)
$ ./obj_dir/rb +verilator+rand+reset+2 +verilator+seed+99
sum_o after reset + 4 clocks = 11111111 (255)
\end{txtcode}

The result now depends on the seed, and \emph{that} is the signal you can
automate against.

\begin{tipbox}{Catch missing resets without four-state simulation}
Run every regression twice with \code{+verilator+rand+reset+2} and two
different \code{+verilator+seed+} values, and require the two runs to produce
byte-identical output. Any dependence on the seed is, by construction, a
dependence on an uninitialised value. This costs one extra run and finds most
of what X-propagation would have found.
\end{tipbox}

\index{high impedance}The other missing state is \code{Z}. Verilator supports
tri-state modelling only in limited forms and will usually report
\code{UNSUPPORTED} or convert the net to a plain wire. Bidirectional pads and
open-drain buses are not its territory.

\subsection{No timing, unless you ask for it}
\label{sec:verilator:timing}

\index{delay!\#}In \VHDL{} you write \code{wait for 5 ns}; in \SV{} you write
\code{\#5}. Verilator's default mode has no notion of time passing inside the
model at all. Give it a delay without saying what to do about it and it
refuses to proceed:

\begin{txtcode}
$ verilator --cc --main --exe --build --top-module delay delay.sv
%Error-NEEDTIMINGOPT: delay.sv:4:5: Use --timing or --no-timing to specify
                                    how delays should be handled
    4 |     #10 a = 1'b1;
      |     ^
%Error: Exiting due to 1 error(s)
\end{txtcode}

With \code{-{}-no-timing} the delay is silently dropped, and Verilator tells
you so:

\begin{txtcode}
%Warning-STMTDLY: delay.sv:4:5: Ignoring delay on this statement due to
                                --no-timing
\end{txtcode}

With \code{-{}-timing}, Verilator compiles delays, \code{wait} statements and
event controls into C++20 coroutines and the program does the right thing:

\begin{txtcode}
$ verilator --binary --timing --top-module delay delay.sv -o dl
$ ./obj_dir/dl
a = 1 at 10
- delay.sv:6: Verilog $finish
\end{txtcode}

\begin{notebox}{\code{-{}-binary} implies \code{-{}-timing}}
In 5.020, \code{-{}-binary} is shorthand for
\code{-{}-main -{}-exe -{}-build -{}-timing}: it generates a \code{main()} for
you and produces an executable in one step. You can see the coroutine support
in the compiler command line it generates, which carries \code{-fcoroutines}
and links \file{verilated\_timing.o}.
\end{notebox}

\index{SDF}\index{gate-level simulation}What this means for gate-level work is
blunt: Verilator does not do it. A \code{specify} block is parsed and
\emph{silently ignored}; there is no SDF back-annotation, no pin-to-pin delay,
no timing check, no \code{\$setuphold}. Post-layout timing simulation stays in
the commercial tool. Verilator's answer to gate-level verification is a
different one: verify the \RTL{} exhaustively and fast, and trust static timing
analysis for the rest -- which is, in practice, what industry does anyway.

\subsection{How much of SystemVerilog does 5.020 actually implement?}
\label{sec:verilator:coverage}

Verilator started as a synthesis-subset tool and has been growing a
verification subset for years. The claims below were each checked on this
machine with a small test file; the error and warning text is quoted verbatim.

\paragraph{Classes, mailboxes, semaphores, fork/join: yes.} The following file
compiles and runs under \code{-{}-binary -{}-timing}:

\begin{svcode}[caption={Class, mailbox, semaphore and fork/join under Verilator 5.020},label={lst:verilator:feat}]
module feat;
  class packet;
    rand bit [7:0] addr;
    rand bit [7:0] data;
    constraint c_addr { addr inside {[16:31]}; }
    function string show();
      return $sformatf("addr=%0d data=%0d", addr, data);
    endfunction
  endclass

  mailbox #(packet) mbx = new(4);
  semaphore         sem = new(1);

  task automatic producer();
    repeat (3) begin
      packet p = new();
      void'(p.randomize());
      mbx.put(p);
      #10;
    end
  endtask

  task automatic consumer();
    repeat (3) begin
      packet p;
      mbx.get(p);
      sem.get(1);
      $display("[%0t] consumed %s", $time, p.show());
      sem.put(1);
    end
  endtask

  initial begin
    fork producer(); consumer(); join
    $display("fork/join completed at %0t", $time);
    $finish;
  end
endmodule
\end{svcode}

\paragraph{Constrained random: no.} The first attempt to build
\cref{lst:verilator:feat} fails, and the message is the important part:

\begin{txtcode}
%Warning-CONSTRAINTIGN: feat.sv:5:16: Constraint ignored (unsupported)
                                    : ... note: In instance 'feat'
    5 |     constraint c_addr { addr inside {[16:31]}; }
      |                ^~~~~~
                     ... For warning description see
                         https://verilator.org/warn/CONSTRAINTIGN?v=5.020
%Error: Exiting due to 1 warning(s)
\end{txtcode}

Add \code{-Wno-CONSTRAINTIGN} and it builds; \code{randomize()} then returns
values that ignore the constraint completely:

\begin{txtcode}
$ ./obj_dir/feat
[0] consumed addr=196 data=156
[10] consumed addr=156 data=2
[20] consumed addr=2 data=228
fork/join completed at 30
- feat.sv:37: Verilog $finish
\end{txtcode}

\code{addr} was constrained to $[16,31]$ and came out as 196. There is no
constraint solver in 5.020: \code{randomize()} is a fancy way of calling
\code{\$urandom} on every \kw{rand} field.

\begin{gotcha}{\code{randomize()} succeeds and lies}
It returns 1. It does not warn at run time. If your testbench relies on
constraints for correctness -- an address that must be aligned, an opcode that
must be legal -- Verilator will feed the \DUT{} illegal stimulus and your
scoreboard will report failures that are not design bugs. Guard the class with
\code{`ifdef VERILATOR} and provide a hand-written randomiser, exactly as
\file{tb\_cordic\_pkg.sv} does with its \code{randomize\_fallback()}.
\end{gotcha}

\paragraph{Covergroups: no.} Not partially, not with warnings -- the parser
rejects them:

\begin{txtcode}
%Error-UNSUPPORTED: cov.sv:3:3: Unsupported: covergroup
    3 |   covergroup cg;
      |   ^~~~~~~~~~
%Error-UNSUPPORTED: cov.sv:4:36: Unsupported: cover bin specification
%Error-UNSUPPORTED: cov.sv:4:12: Unsupported: cover point
\end{txtcode}

Functional coverage in Verilator is done with \code{cover property}, which
\emph{is} supported (\cref{sec:verilator:cov}), or by counting in C++.

\paragraph{Virtual interfaces: yes, but not with clocking blocks.} A plain
virtual interface works:

\begin{svcode}[caption={A virtual interface driven from a class},label={lst:verilator:vif}]
interface simple_if(input logic clk);
  logic [7:0] d;
endinterface

class drv;
  virtual simple_if vif;
  function new(virtual simple_if v); vif = v; endfunction
  task automatic run();
    for (int i = 1; i <= 3; i++) begin
      @(posedge vif.clk);
      vif.d <= 8'(i);
    end
  endtask
endclass
\end{svcode}

\begin{txtcode}
$ ./obj_dir/vif2
virtual interface OK: d=3 at 35
\end{txtcode}

Add a clocking block and hand it out through a modport, as
\file{cordic\_if.sv} does, and 5.020 stops:

\begin{txtcode}
%Error-UNSUPPORTED: cordic_if.sv:70:16: Unsupported: Modport clocking
   70 |   modport drv (clocking drv_cb, output rst_n, input clk);
      |                ^~~~~~~~
\end{txtcode}

Remove the modport and wait directly on the clocking block through the virtual
interface, and you hit an internal error rather than a clean message:

\begin{txtcode}
%Error: Internal Error: vif.sv:13:13: ../V3Width.cpp:2751:
        MemberSel of non-variable
\end{txtcode}

\begin{gotcha}{Clocking blocks do not survive a virtual interface}
The idiom \code{@(vif.drv\_cb); vif.drv\_cb.sig <= x;} is the backbone of a
modern \SV{} testbench and of UVM drivers. In 5.020 it does not work through a
\kw{virtual} interface. You must either drive the raw signals with non-blocking
assignments at \code{@(posedge vif.clk)}, or move the driver into C++ where the
same discipline is enforced by the structure of the evaluation loop. Both
routes are shown later in this chapter.
\end{gotcha}

\paragraph{UVM: no.} UVM needs the constraint solver, covergroups,
\code{uvm\_config\_db} built on parameterised static class members, deep
\code{\$cast}, and a working \code{uvm\_factory}. Several of those are absent.
There is no supported way to run UVM on Verilator 5.020, and you should treat
any claim otherwise with suspicion. \cref{sec:verilator:cppuvm} builds the
missing structure in C++ instead.

\paragraph{Assertions: a useful subset.} Immediate assertions work.
Concurrent assertions with \code{|->}, \code{|=>}, \code{disable iff},
\code{\$stable}, \code{\$past}, \code{\$rose} and \code{\$fell} work. Sequence
operators do not:

\begin{txtcode}
%Error-UNSUPPORTED: cordic_checker.sv:13:56: Unsupported:
        ## () cycle delay range expression
%Error-UNSUPPORTED: cordic_checker.sv:13:58: Unsupported:
        ## (in sequence expression)
\end{txtcode}

\paragraph{A parsing trap worth knowing.} Two constructs in the book's own
\file{tb\_cordic\_pkg.sv} stop the Verilator front end, and both are surprises:

\begin{txtcode}
%Error: tb_cordic_pkg.sv:81:5: Unknown verilator comment:
        '/*verilator does not implement the constraint solver, so the same
        class*/'
%Error: tb_cordic_pkg.sv:283:18: syntax error, unexpected STRING,
        expecting ')' or ','
\end{txtcode}

The first is a \emph{comment}. Any comment whose first word is
\code{verilator} is treated as a metacomment, so a sentence that happens to
begin ``Verilator does not implement\ldots'' becomes a syntax error. The second
is C-style adjacent string literal concatenation inside \code{\$error(\ldots)}.
Here Verilator is right and the source was wrong: 1800-2017 has no rule that
joins two adjacent string literals, and no conforming tool accepts it. The
\SV{} spelling is the concatenation operator, \code{\{"a", "b"\}}, and that is
what \file{tb\_cordic\_pkg.sv} now uses.

\begin{gotcha}{Never begin a comment with the word ``Verilator''}
\code{// Verilator does not support X} is a pragma, not a comment. Write
\code{// This tool does not support X}, or put the word anywhere but first.
\end{gotcha}

\subsection{The comparison, item by item}
\label{sec:verilator:table}

\cref{tab:verilator:vs} summarises. Read the right-hand column as a statement
about Verilator 5.020 specifically; the project moves quickly and several rows
have changed in the last three releases.

\begin{table}[htbp]
  \centering
  \caption{\tool{QuestaSim} and \tool{Verilator} compared. The Verilator column
  reflects version 5.020 as installed and tested for this chapter.}
  \label{tab:verilator:vs}
  \small
  \begin{tabular}{@{}lp{4.4cm}p{5.6cm}@{}}
    \toprule
    \textbf{Aspect} & \textbf{QuestaSim} & \textbf{Verilator 5.020} \\
    \midrule
    Licence and cost   & Commercial, node-locked or floating, per seat
                       & LGPLv3 or Artistic 2.0, free, no server \\
    Execution model    & Event-driven kernel with delta cycles
                       & Static schedule compiled to C++ \\
    \VHDL{} support    & Full IEEE 1076-2008
                       & None; \SV{} and Verilog only \\
    Mixed language     & \VHDL{} and \SV{} in one elaboration
                       & Not possible \\
    Value system       & Four-state (and nine-state for \code{std\_logic})
                       & Two-state; \code{X} becomes 0, 1 or random \\
    Timing and delays  & Full, including transport and inertial
                       & Only with \code{-{}-timing}; ignored otherwise \\
    Gate level and SDF & Full back-annotation and timing checks
                       & Not supported; \code{specify} silently ignored \\
    SVA                & Complete, including sequences and local variables
                       & Useful subset; no \code{\#\#} sequence delay \\
    Covergroups        & Complete
                       & Rejected as \code{UNSUPPORTED} \\
    Constrained random & Full solver with \code{dist}, \code{solve\ldots before}
                       & No solver; \code{CONSTRAINTIGN}, values unconstrained \\
    UVM                & Supported and shipped
                       & Not usable \\
    Classes, mailboxes & Complete
                       & Supported, including \code{fork}/\code{join*} \\
    Virtual interfaces & Complete
                       & Yes, but not with clocking blocks \\
    DPI-C              & Complete
                       & Supported, import and export \\
    VPI                & Complete
                       & Partial; enough for cocotb (\cref{ch:python}) \\
    Waveforms          & WLF, VCD, FSDB
                       & VCD and FST \\
    Interactive debug  & \code{vsim} prompt, breakpoints, force, GUI
                       & None; use \code{gdb} on your own executable \\
    Compile time       & Seconds; incremental per file
                       & Verilation fast, C++ build dominates \\
    Run time           & Baseline
                       & Typically one to two orders of magnitude faster \\
    Memory             & Large per-signal structures
                       & Compact; packed C integers \\
    Multithreading     & Limited
                       & \code{-{}-threads N} partitions the schedule \\
    Coverage           & Line, branch, expression, toggle, functional, GUI
                       & Line, toggle, user; text and annotation only \\
    Continuous integration & Needs a licence server in the runner
                       & A container and a \code{make} target \\
    \bottomrule
  \end{tabular}
\end{table}

\subsection{Why it is fast, and where it is not}
\label{sec:verilator:perf}

\index{Verilator!performance}The speed comes from removing indirection, not
from cleverness at run time:

\begin{enumerate}
  \item There is no event queue, so a signal update is a store to a struct
        member, not an allocation and an insertion into a priority structure.
  \item There is no per-signal sensitivity lookup, because the order of
        evaluation was decided at compile time.
  \item The generated code is compiled by an optimising C++ compiler, which
        inlines, vectorises and keeps hot values in registers.
  \item Activity flags mean that quiet regions cost one branch.
  \item Two-state values pack into machine words, so a 20-bit datapath is one
        \code{uint32\_t} operation rather than an array of four-state cells.
\end{enumerate}

Numbers, measured for this chapter on a two-core Intel Xeon at 2.80\,GHz with
7\,GiB of RAM, Verilator 5.020, \tool{g++} 13.3, generated code compiled at
\code{-O2}. Treat them as an order of magnitude, not as a benchmark: the CORDIC
is a tiny design and your mileage will differ by a large factor in both
directions.

\begin{txtcode}
verilation only (verilator --cc --exe --trace, 3 runs)   0.08 s each
C++ build of the model + testbench (make -j2, -O2)       5.50 s
C++ testbench, 200 000 transactions = 3 200 006 cycles   0.47 s
                                                    ->  6.7 Mcycle/s
SystemVerilog testbench, --binary --timing, same load    6.20 s
                                                    ->  0.54 Mcycle/s
\end{txtcode}

Three lessons are hidden in those five lines.

First, \emph{verilation is not the slow part}. Eighty milliseconds to read the
CORDIC and write the C++; five and a half seconds for \tool{g++} to compile it.
On a large design the ratio stays similar, which is why \code{-j} on the
\code{-{}-build} step matters and why people cache \file{obj\_dir}.

Second, \emph{compile time dominates for small work}. If your regression runs
for two seconds, a five-second build makes Verilator slower end to end than
starting \code{vsim} on an already-compiled library. Verilator wins when the
simulation is long: a processor booting Linux, a DSP chain processing a million
samples, a week-long random regression.

Third, \emph{\code{-{}-timing} costs about an order of magnitude}. The same
\DUT{}, the same number of cycles, driven from a \SV{} class-based testbench
with coroutines rather than from C++, ran at 0.54 instead of 6.7 million cycles
per second. Coroutine suspension and resumption is cheap compared with an event
queue and expensive compared with a function call.

Waveform tracing has a cost of its own. Twenty thousand transactions, about
320\,000 cycles, with every signal traced:

\begin{txtcode}
no tracing                    0.17 s
VCD  (--trace)                0.23 s   106 MB
FST  (--trace-fst)            0.90 s    13 MB
\end{txtcode}

\begin{ruleofthumb}
Trace nothing by default. Add a command-line flag that switches tracing on, and
use FST when the run is long enough that 100\,MB of VCD is a problem.
\end{ruleofthumb}

\section{The Flow, Hands On}
\label{sec:verilator:flow}

\subsection{The switches you will actually type}

\index{Verilator!command line}Three output modes matter.

\begin{itemize}
  \item \code{-{}-cc} emits C++ and nothing else. You supply \code{main()}.
        This is the mode for a C++ testbench and the one this chapter uses
        most.
  \item \code{-{}-binary} emits C++, generates a \code{main()}, enables
        \code{-{}-timing} and builds an executable. This is the mode for
        running a \SV{} testbench.
  \item \code{-{}-sc} emits a SystemC module: the ports become
        \code{sc\_in<>} and \code{sc\_out<>} and the class derives from
        \code{sc\_module}. Useful when the \DUT{} must sit inside an existing
        SystemC virtual platform.
\end{itemize}

Around them:

\begin{itemize}
  \item \code{-{}-top-module <name>} names the top. Verilator can usually
        infer it, but says so noisily when there is more than one candidate;
        always pass it.
  \item \code{-{}-exe file.cpp} says ``this C++ file is part of the program'',
        and makes the generated makefile link an executable.
  \item \code{-Mdir <dir>} changes the output directory from \file{obj\_dir}.
        Use it when you build several variants of the same design, as the
        chapter Makefile does.
  \item \code{-{}-build} runs \tool{make} for you; \code{-j N} parallelises it.
  \item \code{-CFLAGS "\ldots"} and \code{-LDFLAGS "\ldots"} are passed to the
        C++ compiler and linker.
  \item \code{-o <name>} names the executable.
\end{itemize}

\subsection{Verilating the CORDIC}

Assume the sources of \refexamplebook have been copied into \file{rtl/}. The
first command produces C++ only:

\begin{shcode}
$ verilator --cc --top-module cordic_rot rtl/cordic_pkg.sv rtl/cordic_rot.sv
$ ls obj_dir/
Vcordic_rot.cpp    Vcordic_rot__Syms.cpp   Vcordic_rot___024root.h
Vcordic_rot.h      Vcordic_rot__Syms.h     Vcordic_rot___024root__Slow.cpp
Vcordic_rot.mk     Vcordic_rot__pch.h      Vcordic_rot_cordic_pkg.h
Vcordic_rot_classes.mk                     Vcordic_rot_cordic_pkg__Slow.cpp
\end{shcode}

The naming is regular and worth decoding once.
\file{V<top>.h} and \file{V<top>.cpp} are the public wrapper: the class you
instantiate. \file{V<top>\_\_\_024root} is the flattened design; \code{\_\_024}
is the escaped form of \code{\$}, because Verilator names the root scope
\code{\$root}. Files ending in \code{\_\_Slow} contain code that runs once
(construction, reset of internal state, coverage dump) and are compiled with
\code{-Os} rather than \code{-O2}; \code{DepSet} in a file name marks a group
of statements that share a dependency set. \file{V<top>.mk} is the makefile
that \code{-{}-build} invokes.

\subsection{A forty-line testbench}

\cref{lst:verilator:first} is the smallest program that does something useful:
it releases reset, presents one request to the CORDIC, waits for the response
and prints it. Read \code{tick()}-free code once, so that the structure of the
loop is visible before it is abstracted away.

\begin{cppcode}[caption={\file{main\_first.cpp}: the smallest useful Verilator testbench},label={lst:verilator:first}]
#include "Vcordic_rot.h"
#include "verilated.h"
#include <cstdio>

int main(int argc, char** argv) {
    VerilatedContext ctx;
    ctx.commandArgs(argc, argv);
    Vcordic_rot dut{&ctx};

    dut.clk_i = 0;
    dut.rst_ni = 0;
    dut.req_valid_i = 0;
    dut.rsp_ready_i = 1;
    dut.req_x_i = 9949;      // X0_UNIT: seed for cos/sin
    dut.req_y_i = 0;
    dut.req_theta_i = 8579;  // ~pi/3 rad in Q2.13

    for (int cyc = 0; cyc < 200; ++cyc) {
        if (cyc == 4) dut.rst_ni = 1;      // release reset
        if (cyc == 6) dut.req_valid_i = 1; // present one request

        dut.clk_i = 0;                     // low phase: settle the logic
        dut.eval();
        ctx.timeInc(5000);                 // 5 ns, in 1 ps precision units

        bool req_hs = dut.req_valid_i && dut.req_ready_o;
        bool rsp_hs = dut.rsp_valid_o && dut.rsp_ready_i;
        int16_t x = (int16_t)dut.rsp_x_o, y = (int16_t)dut.rsp_y_o;

        dut.clk_i = 1;                     // rising edge
        dut.eval();
        ctx.timeInc(5000);

        if (req_hs) dut.req_valid_i = 0;
        if (rsp_hs) {
            printf("t=%3lu cyc=%d  rsp_x=%6d  rsp_y=%6d\n",
                   (unsigned long)ctx.time(), cyc, x, y);
            break;
        }
    }
    dut.final();
    return 0;
}
\end{cppcode}

\begin{shcode}
$ verilator --cc --exe --build -j 2 --top-module cordic_rot \
      rtl/cordic_pkg.sv rtl/cordic_rot.sv tb/main_first.cpp
$ ./obj_dir/Vcordic_rot
t=220000 cyc=21  rsp_x=  8190  rsp_y= 14190
\end{shcode}

The time stamp is in picoseconds, because \code{timeInc()} counts
\code{timeprecision} units and \file{cordic\_rot.sv} declares
\code{`timescale 1ns/1ps}; 220000\,ps is 220\,ns, that is 22 half periods of
5\,ns.

Seeded with $x_0 = \mathrm{X0\_UNIT} = 9949$ and $y_0 = 0$, the CORDIC returns
$K \cdot 9949 \cdot (\cos\theta, \sin\theta)$, which is
$(\cos\theta, \sin\theta)$ scaled by very nearly $2^{14}$, because
$\mathrm{X0\_UNIT}$ was chosen as $\mathrm{round}(2^{14}/K)$. For $\theta =
8579/2^{13} = 1.0472$\,rad the exact values are $K \cdot 9949 \cdot \cos\theta
= 8191.2$ and $K \cdot 9949 \cdot \sin\theta = 14189.0$. The design is right to
within two LSB, and it took twenty-one clock cycles: four of reset, two idle,
one to accept, fourteen micro-rotations.

\section{The Generated C++ API}
\label{sec:verilator:api}

\subsection{Context, model, ports}

\index{VerilatedContext}\code{VerilatedContext} holds everything that a
simulator kernel would normally own: the current time, the command-line
arguments, the random seed, the coverage database, the ``tracing was ever
enabled'' flag. You may create several, one per model, which is how you run two
independent designs in one process.

\begin{cppcode}
VerilatedContext ctx;
ctx.commandArgs(argc, argv);     // parses +verilator+... plusargs
ctx.traceEverOn(true);           // must precede any trace() call
ctx.timeInc(5000);               // advance time by 5000 precision units
uint64_t now = ctx.time();
Vcordic_rot dut{&ctx};           // the model, bound to this context
\end{cppcode}

If you omit the context, the model attaches to a global default one. That is
fine for a single model and a nuisance for anything else; pass it explicitly.

\index{CData}\index{SData}\index{IData}\index{QData}\index{WData}Ports become
public data members whose C type is chosen by width. The rule is fixed and
worth memorising, because it is the first thing you get wrong:

\begin{table}[htbp]
  \centering
  \caption{Verilator's C types for packed signals.}
  \label{tab:verilator:ctypes}
  \small
  \begin{tabular}{@{}llll@{}}
    \toprule
    \textbf{Width} & \textbf{Verilator type} & \textbf{C type} &
    \textbf{Declared as} \\
    \midrule
    1 to 8    & \code{CData} & \code{uint8\_t}  & \code{VL\_IN8}  \\
    9 to 16   & \code{SData} & \code{uint16\_t} & \code{VL\_IN16} \\
    17 to 32  & \code{IData} & \code{uint32\_t} & \code{VL\_IN}   \\
    33 to 64  & \code{QData} & \code{uint64\_t} & \code{VL\_IN64} \\
    65 and up & \code{WData} & array of \code{uint32\_t} & \code{VL\_INW} \\
    \bottomrule
  \end{tabular}
\end{table}

A module with ports of every width generates exactly this:

\begin{txtcode}
    VL_IN8(&a1,0,0);
    VL_IN8(&a8,7,0);
    VL_IN16(&a16,15,0);
    VL_IN(&a32,31,0);
    VL_INW(&a100,99,0,4);
    VL_OUTW(&y128,127,0,4);
    VL_IN64(&a64,63,0);
\end{txtcode}

\begin{gotcha}{Every port is unsigned, whatever the \SV{} said}
\code{rsp\_x\_o} is declared \code{data\_t}, that is \code{logic signed
[15:0]}. In C++ it is an \code{SData}, that is \code{uint16\_t}. Reading it
without a cast gives you the two's complement bit pattern reinterpreted as a
positive number:
\begin{txtcode}
rsp_y_o as SData  : 51347
rsp_y_o as int16_t: -14189
\end{txtcode}
Always cast a signed output to the matching signed C type, and always mask a
signed input to its width when you write it. The \SV{} sign information does
not cross the boundary.
\end{gotcha}

For a width that is not a multiple of eight, the unused upper bits of the C
type are zero on outputs, and Verilator requires them to be zero on inputs. A
20-bit input written as \code{0xFFFFFFFF} is undefined behaviour, not a
truncation.

\subsection{The evaluation model}

\index{eval()@\code{eval()}}This is where a \VHDL{} intuition fails hardest, so
it gets a numbered recipe.

\code{eval()} does one thing: it brings the model to a settled state,
consistent with the current values of the input members. It is not ``advance
one cycle''. It is not ``advance one delta''. There is no clock inside the
model. \emph{You} are the clock: you write \code{dut.clk\_i = 1} and then call
\code{eval()}, and the fact that \code{clk\_i} went from 0 to 1 between two
calls is what makes the \code{always\_ff} blocks fire.

\begin{enumerate}
  \item Write every input member that should change this cycle.
  \item Call \code{eval()}. Combinational outputs are now valid for the
        \emph{current} values of the state registers.
  \item Read whatever you need to sample \emph{before} the edge. This is the
        C++ equivalent of a clocking block's \code{default input \#1step}.
  \item Set \code{clk\_i = 1}.
  \item Call \code{eval()} again. This is the rising edge: the
        \code{always\_ff} blocks run and the state registers take their new
        values.
  \item Advance the context clock with \code{timeInc()} and, if tracing,
        call \code{dump()} at each of the two time points.
\end{enumerate}

\cref{fig:verilator:eval} draws one cycle of this.

\begin{figure}[htbp]
  \centering
  \begin{tikzpicture}[x=1.55cm,y=1cm]
    % clock waveform
    \draw[thick,draw=codekey]
      (0,0) -- (1,0) -- (1,0.6) -- (2.6,0.6) -- (2.6,0) -- (4,0);
    \node[left=2mm of {(0,0.3)},font=\small\sffamily,anchor=east]
      at (-0.05,0.3) {\code{clk\_i}};

    % time axis
    \draw[-{Latex[length=2mm]},draw=codenumber] (-0.05,-2.6) -- (4.1,-2.6)
      node[right,font=\scriptsize\sffamily]{time};

    % markers
    \foreach \x/\lbl in {0.15/1, 0.55/2, 0.8/3, 1.0/4, 1.25/5}{
      \draw[dashed,draw=codenumber!70] (\x,0.9) -- (\x,-2.5);
    }

    \node[font=\scriptsize\sffamily,anchor=west,align=left]
      at (0.18,-0.5)  {\textbf{1} write inputs};
    \node[font=\scriptsize\sffamily,anchor=west,align=left]
      at (0.58,-1.0)  {\textbf{2} \code{eval()} settles combinational logic};
    \node[font=\scriptsize\sffamily,anchor=west,align=left]
      at (0.83,-1.5)  {\textbf{3} sample outputs (= clocking block input)};
    \node[font=\scriptsize\sffamily,anchor=west,align=left]
      at (1.03,-2.0)  {\textbf{4} \code{clk\_i=1}\quad\textbf{5}
                       \code{eval()}: the edge};

    \node[font=\scriptsize\sffamily,anchor=west]
      at (2.7,-0.5) {\textbf{6} \code{timeInc()}, then the next cycle};

    \node[draw=warnframe,fill=warnbg,rounded corners=2pt,
          font=\scriptsize\sffamily,align=center,inner sep=4pt]
      at (2.05,0.6+1.05)
      {values read \emph{after} step 5 are the \emph{new} register values};
  \end{tikzpicture}
  \caption{One clock cycle in a Verilator C++ testbench. There is no kernel:
  the six steps are six statements you write.}
  \label{fig:verilator:eval}
\end{figure}

\begin{gotcha}{\code{eval()} is not ``advance the clock''}
The commonest beginner's bug is a loop that toggles the clock but calls
\code{eval()} only once per period, or that reads an output after the rising
edge and treats it as the value that existed before it. In \VHDL{} the
distinction is made for you by \code{wait until rising\_edge(clk)} and by the
postponed update of signals. Here it is made by \emph{where in your C++ you put
the read}. Read handshake signals between steps 2 and 4; read registered data
either there or after step 5, but be deliberate about which.
\end{gotcha}

\begin{gotcha}{\code{timeInc()} counts in timeprecision units}
\file{cordic\_rot.sv} declares \code{`timescale 1ns/1ps}. A half period of
5\,ns is therefore \code{ctx.timeInc(5000)}, not \code{ctx.timeInc(5)}. Get it
wrong and everything still runs, but every message the model prints is stamped
with the wrong time. With \code{timeInc(5)} the \DUT{}'s own assertion reports
\begin{txtcode}
[0] %Error: cordic_rot.sv:227: Assertion failed in TOP.cordic_rot:
    cordic_rot: |theta| > pi (theta = 30000)
\end{txtcode}
and with \code{timeInc(5000)} the same failure reports \code{[65000]}, which is
65\,ns. The VCD timestamps are wrong in the same way.

Every C++ testbench in this chapter therefore advances time in picoseconds,
\code{timeInc(5000)} per half period, and every simulated time printed here is
in picoseconds. If you see \code{timeInc(5)} in a testbench that a colleague
wrote, the design still simulates correctly; only the time stamps lie.
\end{gotcha}

The other methods are quickly described. \code{eval\_step()} and
\code{eval\_end\_step()} split \code{eval()} in two, which matters only when
several Verilated models are stepped together in one time slot: you call
\code{eval\_step()} on all of them, then \code{eval\_end\_step()} on all of
them. \code{final()} runs the \SV{} \kw{final} blocks and flushes internal
state; call it once before the program exits, or coverage and \code{\$display}
output can be lost. \code{eventsPending()} and \code{nextTimeSlot()} exist only
in \code{-{}-timing} builds and let you write your own main loop for a design
that has delays in it.

\subsection{Tracing}

\index{VCD}\index{FST}Tracing must be enabled at verilation time
(\code{-{}-trace} for VCD, \code{-{}-trace-fst} for FST) \emph{and} at run
time. Forgetting the first gives a link error rather than a message about
tracing, which is confusing the first time:

\begin{txtcode}
undefined reference to `VerilatedVcd::close()'
undefined reference to `VerilatedTrace<VerilatedVcd,
                        VerilatedVcdBuffer>::dump(unsigned long)'
\end{txtcode}

The run-time sequence is fixed: \code{traceEverOn(true)} before anything else,
then construct the trace object, then \code{dut.trace(tfp, levels)}, then
\code{tfp->open(name)}. After that, call \code{tfp->dump(time)} at every point
you want a sample -- typically twice per cycle, once for each clock phase --
and \code{tfp->close()} at the end.

\begin{cppcode}
#if VM_TRACE_FST
# include "verilated_fst_c.h"
using TraceC = VerilatedFstC;
static const char* TRACE_FILE = "cordic.fst";
#else
# include "verilated_vcd_c.h"
using TraceC = VerilatedVcdC;
static const char* TRACE_FILE = "cordic.vcd";
#endif
\end{cppcode}

\code{VM\_TRACE}, \code{VM\_TRACE\_FST} and \code{VM\_COVERAGE} are defined by
the generated makefile, so one source file can serve both builds.

\begin{tipbox}{Trace only what you need}
\code{dut.trace(tfp, 99)} traces the whole hierarchy, including every
\kw{localparam} in \file{cordic\_pkg.sv} -- the VCD header for this tiny design
already lists \code{ATAN\_ROM[0..13]} and every format constant. Reduce the
depth, or mark the interesting signals with \code{/* verilator public\_flat */}
and trace a shallow level.
\end{tipbox}

\subsection{The dump file is not the viewer}
\label{sec:verilator:viewers}

\index{GTKWave}\index{Surfer}\index{waveform viewer}Nothing above draws a
waveform. \tool{Verilator} writes a \emph{dump file}, and you open that file in
a separate program. This is the one place where the free flow is visibly
different from the commercial one: \tool{QuestaSim} logs signals into its own
\file{.wlf} database and shows them in its own wave window, so the simulator
and the viewer are the same application. With \tool{Verilator} they are not,
and you choose both halves.

There are two dump formats.

\textbf{VCD} is the value-change dump defined by IEEE 1364. It is plain text,
every tool on earth reads it, and it is enormous: a header listing every traced
signal, then one line for every value change. A few thousand cycles of a small
design is a few megabytes; a million cycles of a real one is a file you cannot
open. Use it for short runs, for interoperability, and for anything you want to
read with a script.

\textbf{FST} is the compressed binary format that came out of the \tool{GTKWave}
project. It is typically an order of magnitude smaller than the equivalent VCD
and it opens in a fraction of the time, because a viewer can seek into it
instead of parsing the whole thing. Use it for anything longer than a few
thousand cycles. The measurements earlier in this chapter show the trade in the
other direction: writing FST costs a little more simulation time than writing
VCD, because the compression happens as you go.

Both formats open in the same two viewers.

\begin{itemize}
  \item \tool{GTKWave} is the long-standing one. It reads VCD, FST and several
        other formats, it is in every distribution's package manager, and its
        interface is dated but complete: grouping, radix changes, saved signal
        lists in a \file{.gtkw} file, and a Tcl interface for scripting.
  \item \tool{Surfer} is the newer one. It reads VCD and FST, the interface is
        modern, it starts quickly on large files, and there is a browser build
        as well as a native one, which is genuinely useful when you want to
        send somebody a waveform rather than a screenshot.
\end{itemize}

\begin{shcode}
gtkwave cordic.fst          # or: gtkwave cordic.vcd
surfer   cordic.fst

# QuestaSim, for comparison: the log and the viewer are the same tool
vsim -view vsim.wlf
\end{shcode}

\begin{tipbox}{Save the signal list, not the screenshot}
Both viewers can save the set of signals you have added, their order, their
grouping and their radixes, and reload it against a new dump. Do that once, and
debugging the same block a week later starts from the picture you built rather
than from an empty pane. In \tool{GTKWave} it is \file{File\,>\,Write Save
File}; the \file{.gtkw} belongs in the repository next to the testbench.
\end{tipbox}

\begin{gotcha}{A gigabyte of waveform is not a debugging aid}
\code{\$dumpvars(0, top)} in \SV{}, or \code{dut.trace(tfp, 99)} in C++, dumps
every signal in every instance below that scope, at every change. On the CORDIC
it produces a file you can read; on anything with a memory in it, it produces
tens of gigabytes and a viewer that will not open. Name the scope you care
about --- \code{\$dumpvars(1, u\_dut)} --- or dump for a window of time by
calling \code{\$dumpoff} and \code{\$dumpon} around the interesting cycles.
The equivalent discipline in \tool{QuestaSim} is to \code{log} a subtree
rather than \code{log -r /*}.
\end{gotcha}

\section{Writing a C++ Testbench Properly}
\label{sec:verilator:cpptb}

The forty-line version has all the ideas and none of the discipline. A usable
testbench adds six things: a \code{tick()} helper so the clock protocol is
written once, a reset sequence, a handshake driver that respects
\code{ready}, a golden model, an error counter that becomes the process exit
status, and tracing behind a command-line flag.

\begin{cppcode}[caption={The core of \file{tb\_cordic.cpp}},label={lst:verilator:tb}]
class CordicTb {
  public:
    VerilatedContext ctx;
    Vcordic_rot*     dut;
    TraceC*          tfp = nullptr;
    uint64_t         cycles = 0;
    unsigned         checked = 0, failed = 0;
    double           worst = 0.0;

    void openTrace(const char* name) {
        ctx.traceEverOn(true);
        tfp = new TraceC;
        dut->trace(tfp, 99);
        tfp->open(name);
    }

    // One clock period.  Inputs written before the call are seen by the
    // rising edge; outputs read after the call are the post-edge values.
    void tick() {
        dut->clk_i = 0; dut->eval();
        if (tfp) tfp->dump(ctx.time());
        ctx.timeInc(5000);            // 5 ns at 1 ps precision

        dut->clk_i = 1; dut->eval();
        if (tfp) tfp->dump(ctx.time());
        ctx.timeInc(5000);
        ++cycles;
    }
    // The combinational outputs of the current cycle, before the edge.
    void settle() { dut->clk_i = 0; dut->eval(); }

    void reset(int n = 5) {
        dut->rst_ni = 0;
        for (int i = 0; i < n; ++i) tick();
        dut->rst_ni = 1;
        tick();
    }

    void check(const Txn& t) {
        double th = double(t.theta) / (1 << QA);
        double xr = double(t.x0)    / (1 << QF);
        double yr = double(t.y0)    / (1 << QF);
        double rx = K_GAIN * (xr * std::cos(th) - yr * std::sin(th));
        double ry = K_GAIN * (xr * std::sin(th) + yr * std::cos(th));
        double ex = std::fabs(double(t.x) - rx * (1 << QF));
        double ey = std::fabs(double(t.y) - ry * (1 << QF));
        if (ex > worst) worst = ex;
        if (ey > worst) worst = ey;
        ++checked;
        if (ex > TOL_LSB || ey > TOL_LSB) {
            ++failed;
            printf("MISMATCH x0=%d y0=%d th=%d : got (%d,%d) "
                   "want (%.1f,%.1f) err %.2f/%.2f LSB\n",
                   t.x0, t.y0, t.theta, t.x, t.y,
                   rx * (1 << QF), ry * (1 << QF), ex, ey);
        }
    }
};
\end{cppcode}

The golden model is deliberately written in real arithmetic, at the level of
the specification, exactly as the \SV{} scoreboard of \refverificationbook is. A model
that reimplements the micro-rotations would reproduce the design's bugs.

The handshake is the part that repays care. \code{req\_ready\_o} must be
sampled \emph{before} the edge that consumes the request, which is why
\code{settle()} exists:

\begin{cppcode}
tb.dut->req_valid_i = 1;
tb.dut->req_x_i     = uint16_t(t.x0);
tb.dut->req_y_i     = uint16_t(t.y0);
tb.dut->req_theta_i = uint16_t(t.theta);
bool accepted;
do {
    tb.settle();                       // ready is now this cycle's value
    accepted = tb.dut->req_ready_o;
    tb.tick();                         // the edge that transfers it
} while (!accepted);
tb.dut->req_valid_i = 0;
\end{cppcode}

\begin{gotcha}{Do not test \code{ready} after the edge}
Writing \code{tick(); if (dut->req\_ready\_o) \ldots} tests the value that
\code{ready} takes \emph{after} the transfer, which for this \DUT{} is 0
because it has moved to \code{ST\_ROTATE}. The handshake then never appears to
complete, \code{req\_valid\_i} stays high, and the \DUT{} accepts the same
request again the next time it is idle. The symptom is a testbench that hangs
or that silently sends every transaction twice.
\end{gotcha}

Running it:

\begin{txtcode}
$ verilator --cc --exe --build -j 2 --trace -CFLAGS -O2 \
      --top-module cordic_rot -Mdir obj_tb \
      rtl/cordic_pkg.sv rtl/cordic_rot.sv tb/tb_cordic.cpp -o tb_cordic
$ ./obj_tb/tb_cordic +n=200
----------------------------------------------
 200 transactions, 0 failures, worst error 2.68 LSB
 3206 clock cycles simulated
 *** TEST PASSED ***
$ echo $?
0
\end{txtcode}

\begin{ruleofthumb}
The exit status is the interface to continuous integration. Return non-zero
when anything failed, and never print ``PASSED'' from a run whose exit status
is not zero.
\end{ruleofthumb}

\section{Running a SystemVerilog Testbench: \code{-{}-binary} and \code{-{}-timing}}
\label{sec:verilator:binary}

You do not have to abandon \SV{} for the testbench. With \code{-{}-binary
-{}-timing}, Verilator compiles \kw{initial} blocks, \code{\#} delays,
\code{fork}/\code{join*}, mailboxes and classes into coroutines and generates
its own \code{main()}. The class-based testbench of \refverificationbook is close to
running as it stands. Close, but not quite.

\subsection{What has to change}

The book's \file{tb\_cordic\_pkg.sv} and \file{cordic\_if.sv} need four edits
before 5.020 accepts them, and each one teaches something:

\begin{enumerate}
  \item \textbf{The modport that hands out a clocking block must go}
        (\code{Unsupported: Modport clocking}), and with it every
        \code{@(vif.drv\_cb)} and \code{vif.drv\_cb.sig <= x}. Replace the
        interface with one that carries plain signals, wait on
        \code{@(posedge vif.clk)} and drive with \code{<=}. You lose the
        race-freedom guarantee that the clocking block gave you, so the
        non-blocking assignment now matters: it is what keeps the driver's
        writes in the NBA region, behind the \DUT{}'s sampling.
  \item \textbf{The comment beginning with the word ``Verilator'' must be
        reworded}, because it is parsed as a metacomment.
  \item \textbf{Adjacent string literals in \code{\$error} must be joined}
        into one string.
  \item \textbf{\code{repeat (expr) @(\ldots)} inside a
        \code{fork \ldots join\_none} must become an explicit loop}, because
        of a lifetime restriction on the compiler-generated counter:
\begin{txtcode}
%Error-LIFETIME: tb_vlt_pkg.sv:62:13: Invalid reference:
        Process might outlive variable `__Vrepeat1'.
\end{txtcode}
        The fix is three lines:
\begin{svcode}
int unsigned stall = $urandom_range(3) + 1;
vif.rsp_ready <= 1'b0;
for (int unsigned k = 0; k < stall; k++) @(posedge vif.clk);
\end{svcode}
\end{enumerate}

The parts that need no change at all are worth listing too, because they are
the parts a \VHDL{} designer would expect to be hardest: the classes, the
inheritance, the \code{mailbox \#(cordic\_txn)} between generator and driver,
the queue in the monitor, the \code{fork \ldots join\_none} that starts four
concurrent tasks, the real-valued golden model with \code{\$cos} and
\code{\$sin}, and \code{\$value\$plusargs}. All of that runs.

\subsection{Guarding with \code{`ifdef VERILATOR}}

\index{`ifdef VERILATOR@\code{`ifdef VERILATOR}}Verilator predefines the macro
\code{VERILATOR}, so one source can serve both tools. The book's testbench
already uses this for the two features that do not exist:

\begin{svcode}
`ifndef VERILATOR
  constraint c_theta { theta inside {[-PI_EXT : PI_EXT]}; }
`endif

function void do_randomize();
`ifdef VERILATOR
  randomize_fallback();
`else
  if (!this.randomize()) $fatal(1, "randomization failed");
`endif
endfunction
\end{svcode}

The same pattern wraps the covergroup. Keep the guarded regions small and put
the fallback next to the thing it replaces, so that the two cannot drift apart.

\begin{tipbox}{Make the fallback the same distribution, not just legal}
A hand-written fallback that picks $\theta$ uniformly in $[-\pi, \pi]$ is not
the same stimulus as a \code{dist} that puts weight on the quadrant
boundaries. If the constrained version is what you verify with, and the
Verilator version is what you run nightly, write the corner cases into the
fallback explicitly as directed transactions.
\end{tipbox}

\subsection{Running it}

\begin{txtcode}
$ verilator --binary --timing -j 2 -CFLAGS -O2 --top-module tb_vlt_top \
      rtl/cordic_pkg.sv svtb/cordic_vif.sv rtl/cordic_rot.sv \
      svtb/tb_vlt_pkg.sv svtb/tb_vlt_top.sv -o tb_svtb
$ ./obj_svtb/tb_svtb +n=200
 scoreboard: 200 checked, 0 failed
 worst error: 2.82 LSB (tolerance 12 LSB)
 *** TEST PASSED ***
 simulated time: 33525000
- svtb/tb_vlt_top.sv:33: Verilog $finish
\end{txtcode}

The whole layered testbench -- generator, driver with random back-pressure,
passive monitor, scoreboard with a trigonometric reference model -- runs
unmodified in structure on a free tool. That is a genuinely useful result, and
it is worth weighing against the two things you gave up: constrained random
and functional coverage.

\begin{notebox}{When to choose the \SV{} testbench over the C++ one}
Choose \SV{} when the testbench must also run in \tool{QuestaSim} (one source,
two tools), when the protocol is complex enough that blocking tasks are much
clearer than state machines, or when the people who will maintain it are \RTL{}
engineers. Choose C++ when the run is long and speed matters, when the golden
model already exists in C or C++, or when you want to link the \DUT{} into a
larger software system.
\end{notebox}

\section{A UVM-Shaped Verification Environment in C++}
\label{sec:verilator:cppuvm}

UVM is not available. That does not mean the \emph{ideas} in UVM are
unavailable: they are ordinary object-oriented design, and C++ expresses them
with less ceremony than \SV{} does. This section builds a small, complete,
header-only framework and runs it against the CORDIC. It is about 200 lines.

\subsection{What UVM is actually for}

Strip away the macros and the phasing and UVM does four things (\refverificationbook):

\begin{enumerate}
  \item It gives components a \textbf{name and a place in a tree}, and calls
        \code{build\_phase}, \code{connect\_phase} and \code{run\_phase} on
        every one of them in a defined order.
  \item It provides \textbf{TLM ports} so a monitor can publish transactions
        without knowing who consumes them.
  \item It provides a \textbf{factory}, so a test can substitute a different
        implementation of a component without editing the environment.
  \item It provides a \textbf{configuration database}, so a test can change a
        parameter deep in the hierarchy without editing anything.
\end{enumerate}

Points 3 and 4 are the reason UVM exists. Everything else is convenience.
\cref{fig:verilator:env} shows the hierarchy we will build, deliberately drawn
the same way as the UVM hierarchy in \refverificationbook.

\begin{figure}[htbp]
  \centering
  \begin{tikzpicture}[
      every node/.style={font=\small\sffamily},
      comp/.style={draw=codeframe,fill=codebg,rounded corners=2pt,
                   minimum height=7mm,minimum width=22mm,inner xsep=4pt},
      envb/.style={draw=codekey,fill=notebg,rounded corners=2pt,
                   minimum height=7mm,minimum width=22mm},
      agt/.style={draw=accentalt,fill=warnbg,rounded corners=2pt,
                  minimum height=7mm,minimum width=22mm},
      ar/.style={-{Latex[length=1.8mm]},draw=codenumber},
      tlm/.style={-{Latex[length=1.8mm]},draw=accent,thick,dashed}]

    \node[envb] (test) at (0,0) {\code{CordicTest}};
    \node[envb] (env)  at (0,-1.2) {\code{CordicEnv}};
    \node[agt]  (agt)  at (-2.6,-2.6) {\code{CordicAgent}};
    \node[comp] (sb)   at (2.2,-2.6) {\code{CordicSb}};
    \node[comp] (cov)  at (2.2,-3.8) {\code{CordicCov}};

    \node[comp] (seq)  at (-5.0,-4.2) {\code{CordicSeq}};
    \node[comp] (drv)  at (-2.6,-4.2) {\code{CordicDrv}};
    \node[comp] (mon)  at (-0.2,-4.2) {\code{CordicMon}};

    \node[draw=accent,fill=tipbg,rounded corners=2pt,minimum height=7mm]
         (dut) at (-1.4,-5.8) {\code{Vcordic\_rot} (the model)};

    \draw[ar] (test) -- (env);
    \draw[ar] (env)  -- (agt);
    \draw[ar] (env)  -- (sb);
    \draw[ar] (env)  -- (cov);
    \draw[ar] (agt)  -- (seq);
    \draw[ar] (agt)  -- (drv);
    \draw[ar] (agt)  -- (mon);

    \draw[ar,densely dotted] (seq) -- node[below,font=\tiny]{queue} (drv);
    \draw[ar] (drv.south) -- (dut.north west);
    \draw[ar] (dut.north east) -- (mon.south);

    \draw[tlm] (mon.east) .. controls (1.2,-4.2) and (1.2,-2.9) .. (sb.south west);
    \draw[tlm] (mon.east) .. controls (1.4,-4.4) and (1.4,-3.9) .. (cov.west);
    \node[font=\tiny\sffamily,text=accent] at (1.35,-3.25)
      {\code{AnalysisPort<CordicTxn>}};
  \end{tikzpicture}
  \caption{The C++ verification environment. Solid arrows are ownership,
  dashed green arrows are TLM connections. Compare with the UVM hierarchy in
  \refverificationbook: the boxes are the same.}
  \label{fig:verilator:env}
\end{figure}

\subsection{The framework}

\index{TLM}\index{factory}The whole framework is one header. The component base
class carries a name, a parent, a list of children, and four virtual phase
methods.

\begin{cppcode}[caption={\file{svtlm.h}: the component base class},label={lst:verilator:comp}]
class Component {
  public:
    Component(std::string name, Component* parent)
        : name_(std::move(name)), parent_(parent) {
        if (parent_) parent_->children_.push_back(this);
    }
    virtual ~Component() = default;

    std::string fullName() const {
        return parent_ ? parent_->fullName() + "." + name_ : name_;
    }

    // ---- the phases, overridden by the user -------------------------
    virtual void build()   {}
    virtual void connect() {}
    virtual void tick()    {}
    virtual void report()  {}

    // ---- the phase walkers, called by the test ----------------------
    void buildAll() {
        build();
        // build() may have created children, so index by hand.
        for (size_t i = 0; i < children_.size(); ++i) children_[i]->buildAll();
    }
    void connectAll() {                    // bottom-up, like UVM
        for (auto* c : children_) c->connectAll();
        connect();
    }
    void tickAll() {
        tick();
        for (auto* c : children_) c->tickAll();
    }
    void reportAll() {
        for (auto* c : children_) c->reportAll();
        report();
    }
    void info(const char* fmt, ...) const;   // prints "[full.name] ..."
    int  cfgInt(const std::string& field, int dflt) const;

  protected:
    std::vector<Component*> children_;
  private:
    std::string name_;
    Component*  parent_;
};
\end{cppcode}

The analysis port is nine lines, because C++ already has first-class callables.
In UVM you would declare an \code{uvm\_analysis\_imp} with a macro so that
several subscribers could each have a distinct \code{write} method; here a
lambda does it.

\begin{cppcode}[caption={The TLM analysis port},label={lst:verilator:ap}]
template <typename T>
class AnalysisPort {
  public:
    using Sub = std::function<void(const T&)>;
    void connect(Sub s) { subs_.push_back(std::move(s)); }
    void write(const T& t) const { for (const auto& s : subs_) s(t); }
  private:
    std::vector<Sub> subs_;
};
\end{cppcode}

The factory maps a type \emph{name} to a constructor, plus a table of
overrides. Registration is one line in each component, hidden behind a macro
that is the direct analogue of \code{`uvm\_component\_utils}.

\begin{cppcode}[caption={Factory and registration},label={lst:verilator:factory}]
class Factory {
  public:
    using Ctor = std::function<Component*(const std::string&, Component*)>;
    static Factory& get() { static Factory f; return f; }

    void reg(const std::string& type, Ctor c) { ctors_[type] = std::move(c); }
    void setTypeOverride(const std::string& orig, const std::string& repl) {
        overrides_[orig] = repl;
    }
    Component* create(const std::string& type, const std::string& name,
                      Component* parent) {
        std::string t = type;
        auto ov = overrides_.find(t);
        if (ov != overrides_.end()) t = ov->second;
        auto it = ctors_.find(t);
        if (it == ctors_.end()) { printf("FATAL: no type '%s'\n", t.c_str());
                                  abort(); }
        return it->second(name, parent);
    }
    template <typename T>
    T* createAs(const std::string& type, const std::string& name,
                Component* parent) {
        return static_cast<T*>(create(type, name, parent));
    }
  private:
    std::map<std::string, Ctor>        ctors_;
    std::map<std::string, std::string> overrides_;
};

template <typename T>
struct Registrar {
    explicit Registrar(const std::string& type) {
        Factory::get().reg(type, [](const std::string& n, Component* p) {
            return static_cast<Component*>(new T(n, p));
        });
    }
};

#define SVTLM_REGISTER(TYPE)                                     \
    std::string typeName() const override { return #TYPE; }      \
    static inline svtlm::Registrar<TYPE> svtlm_reg_{#TYPE}
\end{cppcode}

The configuration store is a pair of \code{std::map}s keyed by
\code{"<component path>.<field>"}, with \code{"*"} as a wildcard path. A
component reads it during \code{build()} through \code{cfgInt()}, which is
where UVM reads \code{uvm\_config\_db} too.

\subsection{The hard part: blocking tasks versus callbacks}

\index{coroutine}Here is the real difficulty, and it deserves to be stated
plainly rather than designed around silently.

A \SV{} driver is a \emph{blocking task}. It says ``drive these pins, then wait
until \code{ready}, then wait one more cycle, then drop \code{valid}''. The
program counter of that task \emph{is} the state of the protocol. Nothing else
has to remember where the driver had got to.

A C++ method is not blocking. \code{tick()} is called, it returns, and every
piece of state that must survive until the next call has to be stored in a
member variable. There are two ways out.

\paragraph{(a) The tick model.} Every component owns an explicit state machine
and a \code{tick()} that advances it by one clock. This is portable to any C++
compiler, needs no library, and debugs well because the state is a named
variable you can print. It costs you readability: the protocol is no longer a
straight line of code.

\begin{cppcode}[caption={The driver as an explicit state machine},label={lst:verilator:drv}]
class CordicDrv : public Component {
  public:
    SVTLM_REGISTER(CordicDrv);
    CordicDrv(const std::string& n, Component* p) : Component(n, p) {}

    CordicVif*            vif = nullptr;
    std::deque<CordicTxn> q;
    unsigned              stall_pct = 0, n_driven = 0;

    void build() override {
        stall_pct = cfgInt("stall_pct", 0);
        rng_.seed(cfgInt("seed", 2));
    }
    void tick() override {
        Vcordic_rot* d = vif->dut;

        // --- request channel -----------------------------------------
        if (st_ == IDLE) {
            if (!q.empty()) {
                const CordicTxn& t = q.front();
                d->req_valid_i = 1;
                d->req_x_i     = uint16_t(t.x0);
                d->req_y_i     = uint16_t(t.y0);
                d->req_theta_i = uint16_t(t.theta);
                st_ = DRIVING;
            } else {
                d->req_valid_i = 0;
            }
        } else if (d->req_ready_o) {   // handshake completes on this edge
            q.pop_front();
            ++n_driven;
            d->req_valid_i = 0;
            st_ = IDLE;
        }

        // --- response channel: random back-pressure -------------------
        d->rsp_ready_i = (stall_pct == 0)
                       || (int(rng_() % 100) >= int(stall_pct));
    }
    void report() override { info("drove %u transactions", n_driven); }

  private:
    enum State { IDLE, DRIVING } st_ = IDLE;
    std::mt19937 rng_{2};
};
\end{cppcode}

The monitor is written the same way and is where the analysis port lives:

\begin{cppcode}[caption={The monitor publishes on the analysis port},label={lst:verilator:mon}]
void tick() override {
    Vcordic_rot* d = vif->dut;
    if (!d->rst_ni) { pending_.clear(); return; }

    if (d->req_valid_i && d->req_ready_o) {
        CordicTxn t;
        t.x0 = int16_t(d->req_x_i);
        t.y0 = int16_t(d->req_y_i);
        t.theta = int16_t(d->req_theta_i);
        t.id = n_seen++;
        pending_.push_back(t);
    }
    if (d->rsp_valid_o && d->rsp_ready_i) {
        if (pending_.empty()) { info("ERROR response with no request"); return; }
        CordicTxn t = pending_.front();
        pending_.pop_front();
        t.x = int16_t(d->rsp_x_o);
        t.y = int16_t(d->rsp_y_o);
        ap.write(t);                       // TLM broadcast
    }
}
\end{cppcode}

The environment does the wiring, and this is where the TLM idea earns its
keep: the monitor knows nothing about the scoreboard or the coverage
collector.

\begin{cppcode}
void connect() override {
    CordicSb*  s = sb;
    CordicCov* c = cov;
    agt->mon->ap.connect([s](const CordicTxn& t) { s->write(t); });
    agt->mon->ap.connect([c](const CordicTxn& t) { c->write(t); });
}
\end{cppcode}

The clock loop is the only place that knows about both the model and the
components. Note where \code{tickAll()} sits: after the settling
\code{eval()} and before the rising edge, so every component sees exactly the
pre-edge values, which is the semantics of a clocking block with
\code{default input \#1step}.

\begin{cppcode}[caption={The clock loop, in the test},label={lst:verilator:run}]
void run(VerilatedContext& ctx, unsigned max_cycles) {
    Vcordic_rot* d = vif.dut;
    d->clk_i = 0; d->rst_ni = 0;
    d->req_valid_i = 0; d->rsp_ready_i = 0;
    d->req_x_i = d->req_y_i = d->req_theta_i = 0;

    for (unsigned c = 0; c < max_cycles; ++c) {
        d->clk_i = 0;
        d->eval();                 // settle: pre-edge values are now valid
        if (c == 5) d->rst_ni = 1;
        if (c > 5) tickAll();      // every component sees this cycle
        ctx.timeInc(5000);         // 5 ns, in 1 ps precision units
        d->clk_i = 1;
        d->eval();                 // the rising edge
        ctx.timeInc(5000);
        if (finished()) break;
    }
}
\end{cppcode}

\begin{gotcha}{A \code{tick()} that writes an input cannot then read its effect}
Inside \code{tick()} you write \code{req\_valid\_i} and read
\code{req\_ready\_o} in the same breath. The value you read was computed by the
\code{eval()} that ran \emph{before} \code{tickAll()}, so it does not include
your write. For a well-formed valid/ready interface that is correct and
desirable: it is exactly what a clocking block gives you. For a combinational
handshake where \code{ready} depends on \code{valid}, it is wrong, and you must
either call \code{eval()} again inside the tick or split the driver into two
phases. Know which kind of interface you have.
\end{gotcha}

\paragraph{(b) C++20 coroutines.} The blocking style can be recovered. A
coroutine suspends at \code{co\_await} and resumes where it left off, which is
precisely what a \SV{} task does at \code{@(posedge clk)}. Verilator itself
uses this mechanism internally for \code{-{}-timing}, and \tool{g++} 13 on this
machine supports it with \code{-std=c++20 -fcoroutines}. A scheduler is three
lines:

\begin{cppcode}[caption={A three-line cooperative scheduler},label={lst:verilator:coro}]
static std::vector<std::coroutine_handle<>> g_waiting;

struct AwaitClock {                    // == SystemVerilog's @(posedge clk)
    bool await_ready() const noexcept { return false; }
    void await_suspend(std::coroutine_handle<> h) const {
        g_waiting.push_back(h);
    }
    void await_resume() const noexcept {}
};

struct Task {                          // == a SystemVerilog task
    struct promise_type {
        Task get_return_object() noexcept { return {}; }
        std::suspend_never initial_suspend() noexcept { return {}; }
        std::suspend_never final_suspend() noexcept { return {}; }
        void return_void() noexcept {}
        void unhandled_exception() { std::terminate(); }
    };
};

static void resumeAll() {              // one clock edge for every process
    std::vector<std::coroutine_handle<>> ready;
    ready.swap(g_waiting);
    for (auto h : ready) h.resume();
}
\end{cppcode}

The driver then reads almost exactly like the \SV{} task it replaces:

\begin{cppcode}[caption={The same driver, in blocking style},label={lst:verilator:codrv}]
Task driver(Vcordic_rot* d) {          // == task automatic run();
    for (;;) {
        while (g_queue.empty()) co_await AwaitClock{};
        const Txn& t = g_queue.front();
        d->req_valid_i = 1;
        d->req_x_i     = uint16_t(t.x0);
        d->req_y_i     = uint16_t(t.y0);
        d->req_theta_i = uint16_t(t.theta);
        // req_ready_o for THIS cycle is already visible: wait until it is
        // high, then let one more edge pass -- that edge is the transfer.
        while (!d->req_ready_o) co_await AwaitClock{};
        co_await AwaitClock{};
        d->req_valid_i = 0;
        g_queue.pop_front();
    }
}
\end{cppcode}

Starting three such processes is the C++ spelling of
\code{fork \ldots join\_none}:

\begin{cppcode}
driver(&d);
monitor(&d);
sequence(n);
for (; cyc < limit && g_checked < n; ++cyc) {
    d.clk_i = 0; d.eval();
    if (cyc == 5) d.rst_ni = 1;
    if (cyc > 5) resumeAll();          // every process gets its clock edge
    ctx.timeInc(5000);
    d.clk_i = 1; d.eval();
    ctx.timeInc(5000);
}
\end{cppcode}

It compiles and it runs, and it produces the same numbers as the tick version:

\begin{txtcode}
$ verilator --cc --exe --build -j 2 -CFLAGS "-std=c++20 -fcoroutines -O2" \
      --top-module cordic_rot -Mdir obj_coro \
      rtl/cordic_pkg.sv rtl/cordic_rot.sv coro/main_coro.cpp -o tb_coro
$ ./obj_coro/tb_coro +n=200
coroutine TB: 200 checked, 0 failed, worst 2.68 LSB, 3206 cycles
*** TEST PASSED ***
\end{txtcode}

Be honest about the cost. The forty lines above are the easy part; a real
framework needs \code{join\_any}, \code{join\_none} with a handle you can kill,
\code{wait} on an arbitrary condition, timeouts, and a way to propagate an
exception out of a suspended coroutine. Lifetime is the sharp edge: a coroutine
frame is heap-allocated and outlives the function that created it, so a
reference captured before a \code{co\_await} may dangle after it, and the
compiler will not tell you. Measured on the same 200\,000-transaction load, the
coroutine version ran in 0.66\,s against 0.47\,s for the tick version -- about
40\% slower, still an order of magnitude faster than \code{-{}-binary
-{}-timing}.

\begin{ruleofthumb}
Write the tick model first. Reach for coroutines only when a protocol has so
many states that the state machine has stopped being readable.
\end{ruleofthumb}

\subsection{The factory and the config store in use}

The test configures before it builds, exactly as a UVM test does in
\code{new()} or in an early phase:

\begin{cppcode}
ConfigDb& cfg = ConfigDb::get();
cfg.setInt("*",           "n_txn",    n);
cfg.setInt("*",           "stall_pct", stall);
cfg.setInt("test.env.sb", "tol_lsb",  12);

if (strict) Factory::get().setTypeOverride("CordicSb", "CordicSbStrict");

Vcordic_rot dut{&ctx};
CordicTest  test("test", nullptr);
test.vif.dut = &dut;
test.buildAll();
test.connectAll();
test.printTree();
\end{cppcode}

Running with default settings:

\begin{txtcode}
$ ./obj_uvm/tb_uvm +n=200 +stall=25
[test.env.agt.seq      ] will generate 200 transactions
--- component tree ---
test (CordicTest)
  env (CordicEnv)
    agt (CordicAgent)
      seq (CordicSeq)
      drv (CordicDrv)
      mon (CordicMon)
    sb (CordicSb)
    cov (CordicCov)
--- run ---
--- report ---
[test.env.agt.drv      ] drove 200 transactions
[test.env.sb           ] 200 checked, 0 failed, worst 2.68 LSB (tol 12)
[test.env.cov          ] theta bins hit 4/5  [62 47 0 46 45]
*** TEST PASSED ***
\end{txtcode}

Now the same binary with \code{+strict}, which registers a type override so
that the factory builds a \code{CordicSbStrict} -- a subclass with a tolerance
of 2 LSB -- wherever the environment asked for a \code{CordicSb}. Not one line
of \code{CordicEnv} changes:

\begin{txtcode}
$ ./obj_uvm/tb_uvm +n=200 +strict
--- component tree ---
test (CordicTest)
  env (CordicEnv)
    agt (CordicAgent)
      seq (CordicSeq)
      drv (CordicDrv)
      mon (CordicMon)
    sb (CordicSbStrict)
    cov (CordicCov)
--- run ---
[test.env.sb  ] MISMATCH id=33 theta=1276 got(-8109,-9602)
                                          want(-8107.0,-9602.8)
--- report ---
[test.env.agt.drv      ] drove 200 transactions
[test.env.sb           ] STRICT mode: 200 checked, 4 failed,
                         worst 2.68 LSB (tol 2)
*** TEST FAILED ***
$ echo $?
1
\end{txtcode}

That is the whole point of a factory, demonstrated in eleven lines of framework
code. The coverage collector reports \code{theta bins hit 4/5}: the bin for
$\theta = 0$ exactly was never hit, which is the kind of thing functional
coverage exists to tell you.

\subsection{An honest comparison}

\begin{table}[htbp]
  \centering
  \caption{The C++ framework of this chapter against UVM.}
  \label{tab:verilator:vsuvm}
  \small
  \begin{tabular}{@{}lp{4.2cm}p{5.4cm}@{}}
    \toprule
    \textbf{Feature} & \textbf{UVM 1.2 / IEEE 1800.2} & \textbf{This framework} \\
    \midrule
    Size to learn   & About 400 classes and 200 macros
                    & 4 classes, 1 macro, 200 lines \\
    Phasing         & 9 common phases plus 12 run-time sub-phases,
                      with objections
                    & 4 phases, called explicitly by the test \\
    Concurrency     & Native blocking tasks
                    & \code{tick()} state machines, or C++20 coroutines \\
    TLM             & Ports, exports, imps, FIFOs, analysis
                    & \code{AnalysisPort<T>} only \\
    Factory         & Type and instance override, by name and by type
                    & Type override by name \\
    Configuration   & \code{uvm\_config\_db\#(T)}, hierarchical wildcards
                    & Two \code{std::map}s, one wildcard \\
    Sequences       & Sequence library, layering, arbitration, virtual
                      sequences
                    & A component that pushes into a queue \\
    Constrained random & Full solver
                    & Whatever \code{<random>} gives you \\
    Coverage        & Covergroups, cross coverage, database merging
                    & Counters and \code{cover property} \\
    Register layer  & \code{uvm\_reg}, adapters, predictors, built-in tests
                    & Nothing \\
    Reporting       & Severity, verbosity, ID filtering, catchers
                    & \code{printf} with the component path \\
    Reuse across projects & Real, and the industry expects it
                    & Only what you carry yourself \\
    Speed           & Baseline
                    & One to two orders of magnitude faster \\
    Tool support    & Every commercial simulator
                    & Any C++17 compiler \\
    \bottomrule
  \end{tabular}
\end{table}

\begin{notebox}{When this is the right choice}
Use a C++ environment when the reference model is already C or C++ (a DSP
kernel, an instruction-set simulator, a codec), when the regression is long
enough that speed dominates, when there is no licence, or when the design is
small enough that UVM's overhead is larger than the design. Use UVM when the
project is large, when several teams must reuse each other's agents, when a
register model would otherwise be written by hand, or when the customer
requires it. Do not use this framework as a substitute for learning UVM: learn
UVM, then recognise that you can rebuild its useful 10\% in an afternoon.
\end{notebox}

\section{Coverage and Assertions}
\label{sec:verilator:cov}

\subsection{Collecting coverage}

\index{coverage!line}\index{coverage!toggle}Three switches turn coverage on:

\begin{itemize}
  \item \code{-{}-coverage-line} counts how many times each basic block was
        entered.
  \item \code{-{}-coverage-toggle} counts 0-to-1 and 1-to-0 transitions on
        every bit of every signal.
  \item \code{-{}-coverage-user} counts \code{cover property} statements that
        you wrote.
\end{itemize}

\code{-{}-coverage} is shorthand for all three. In a \code{-{}-cc} build the
database is not written automatically: your program must ask for it, which is
one more consequence of ``Verilator is a compiler''.

\begin{cppcode}
#if VM_COVERAGE
    Verilated::mkdir("logs");
    ctx.coveragep()->write("logs/coverage.dat");
#endif
\end{cppcode}

The CORDIC has no cover points of its own, so we bind some onto it. A separate
checker module keeps the \RTL{} clean and is the same technique you would use
in \tool{QuestaSim}:

\begin{svcode}[caption={Assertions and cover points bound onto the DUT},label={lst:verilator:chk}]
module cordic_checker (
  input logic clk_i, rst_ni,
  input logic req_valid_i, req_ready_o,
  input logic rsp_valid_o, rsp_ready_i
);
  // --- functional cover points (--coverage-user) ------------------------
  c_req_accept : cover property (@(posedge clk_i) disable iff (!rst_ni)
                                 req_valid_i && req_ready_o);
  c_backpressure : cover property (@(posedge clk_i) disable iff (!rst_ni)
                                   rsp_valid_o && !rsp_ready_i);

  // --- assertions (--assert) --------------------------------------------
  a_reset_quiet : assert property (@(posedge clk_i) (!rst_ni) |-> !rsp_valid_o)
    else $error("rsp_valid_o high during reset");
  a_no_ready_drop : assert property (@(posedge clk_i) disable iff (!rst_ni)
                                     (rsp_valid_o && !rsp_ready_i) |=> rsp_valid_o)
    else $error("rsp_valid_o dropped before it was accepted");
endmodule

bind cordic_rot cordic_checker u_chk (.*);
\end{svcode}

\code{bind} works, and \code{.*} works with it. Build with coverage on, run,
and annotate:

\begin{txtcode}
$ verilator --cc --exe --build -j 2 --trace --assert \
      --coverage-line --coverage-toggle --coverage-user \
      --top-module cordic_rot -Mdir obj_cov \
      rtl/cordic_pkg.sv rtl/cordic_rot.sv cov/cordic_checker.sv \
      tb/tb_cordic.cpp -o tb_cov
$ ./obj_cov/tb_cov +n=300
$ verilator_coverage --annotate annotated logs/coverage.dat
Total coverage (72/118) 61.00%
See lines with '%00' in annotated
\end{txtcode}

\index{verilator\_coverage@\code{verilator\_coverage}}\code{verilator\_coverage
-{}-annotate} writes a copy of every source file with a count in the left
margin. A line prefixed \code{\%000000} was never executed:

\begin{txtcode}
        module cordic_checker (
%000001   input logic clk_i, rst_ni,
+009611  point: comment=clk_i
-000001  point: comment=rst_ni
 000600   input logic req_valid_i, req_ready_o,
+000600  point: comment=req_valid_i
+000601  point: comment=req_ready_o
        );
 000300   c_req_accept : cover property (@(posedge clk_i) disable iff (!rst_ni)
+000300  point: comment=c_req_accept
%000000   c_backpressure : cover property (@(posedge clk_i) disable iff (!rst_ni)
-000000  point: comment=c_backpressure
\end{txtcode}

Read that carefully, because it is coverage doing its job. \code{c\_req\_accept}
fired 300 times, once per transaction. \code{c\_backpressure} fired zero times:
the testbench holds \code{rsp\_ready\_i} at 1 forever and has therefore never
tested what happens when the consumer stalls. Adding four lines of random
back-pressure to the testbench and re-running:

\begin{txtcode}
$ ./obj_cov/tb_cov +n=300 +bp
$ verilator_coverage --annotate annotated logs/coverage.dat
Total coverage (76/118) 64.00%
$ grep -A1 c_backpressure annotated/cordic_checker.sv
 000120   c_backpressure : cover property (@(posedge clk_i) disable iff (!rst_ni)
 000120                                    rsp_valid_o && !rsp_ready_i);
\end{txtcode}

\begin{gotcha}{Toggle coverage on a reset signal is always 1}
Look at the annotation above: \code{rst\_ni} shows a count of 1, and it is
marked \code{-} rather than \code{+}. Reset toggles once per simulation. That
is not a coverage hole and no amount of stimulus will change it. Learn to
recognise the points that are structurally unreachable before you start chasing
a coverage percentage, and waive them in the \file{.vlt} file.
\end{gotcha}

\subsection{Assertions}

\index{assertion!concurrent}\code{-{}-assert} enables both the assertions you
bound on and the ones the designer wrote inside the module. Without it they are
compiled away entirely. \file{cordic\_rot.sv} contains an immediate assertion
on the input range; here is what happens when the testbench drives
$\theta = 30000$, which is larger than $\pi \cdot 2^{13} = 25736$:

\begin{txtcode}
$ ./obj_as/tb_as              # built with --assert
[65000] %Error: cordic_rot.sv:227: Assertion failed in TOP.cordic_rot:
        cordic_rot: |theta| > pi (theta = 30000)
%Error: rtl/cordic_rot.sv:227: Verilog $stop
Aborting...
$ echo $?
134

$ ./obj_na/tb_na              # built without --assert
$ echo $?
0
\end{txtcode}

\begin{tipbox}{Always build the regression with \code{-{}-assert}}
The cost is small and the benefit is that every assertion the designers wrote
becomes a test. Note also that a failing assertion calls \code{\$stop}, which
aborts the process; if you want the run to continue and count failures, use
\code{\$error} without an enclosing \code{assert}, or catch the abort in your
build script.
\end{tipbox}

The supported subset is the practically useful one: implication (\code{|->} and
\code{|=>}), \code{disable iff}, \code{\$past}, \code{\$stable}, \code{\$rose},
\code{\$fell}, \code{\$onehot} and immediate assertions inside procedural code.
The missing piece is sequences: \code{a \#\#1 b} and \code{a \#\#[1:3] b} are
rejected as \code{UNSUPPORTED} in 5.020, which rules out most multi-cycle
protocol properties. Where you need one, write it as a small checker \RTL{}
module with an explicit state machine and an immediate assertion; that also
works in every other tool.

\section{Linting as a First-Class Use}
\label{sec:verilator:lint}

\index{lint}\index{Verilator!lint}Many teams that never simulate with Verilator
use it every day, because \code{verilator -{}-lint-only -Wall} is the best free
static checker for \SV{} and it takes under a second on a large design.

\refdesignbook argued that \VHDL{}'s value comes largely from its type
system: \code{std\_logic\_vector} does not silently become an \code{integer},
ranges are checked, and an unconstrained port is a compile error. \SV{} gives
almost all of that up. Widths adapt silently, signedness is contagious in ways
that surprise, an undeclared identifier becomes a one-bit wire, and an
incomplete \code{case} produces a latch without comment. \emph{Lint is the
cultural replacement for the \VHDL{} type checker.} A team that writes \SV{}
without a lint gate in continuous integration has thrown away a safety net and
put nothing in its place.

\subsection{The warnings a VHDL designer meets first}

The module in \cref{lst:verilator:sloppy} is a catalogue of first-week
mistakes.

\begin{svcode}[caption={\file{sloppy.sv}: every mistake at once},label={lst:verilator:sloppy}]
module sloppy (
  input  logic        clk,
  input  logic [7:0]  a,
  input  logic [1:0]  sel,
  output logic [3:0]  y,
  output logic [7:0]  q
);
  logic [3:0] narrow;
  logic [9:0] wide;

  assign narrow = a;          // 8 bits into 4
  assign wide   = a;          // 8 bits into 10

  always_comb begin           // incomplete case
    case (sel)
      2'b00: y = 4'd1;
      2'b01: y = 4'd2;
      2'b10: y = 4'd3;
    endcase
  end

  always_ff @(posedge clk) begin
    q = a;                    // blocking in a clocked process
  end

  assign z = a[0];            // implicit net

  logic [3:0] loop;
  always_comb loop = {loop[2:0], narrow[3]};
endmodule
\end{svcode}

\begin{txtcode}
$ verilator --lint-only -Wall --top-module sloppy sloppy.sv
%Warning-IMPLICIT: sloppy.sv:32:10: Signal definition not found,
                   creating implicitly: 'z'
%Warning-WIDTHTRUNC: sloppy.sv:14:17: Operator ASSIGNW expects 4 bits on
                     the Assign RHS, but Assign RHS's VARREF 'a'
                     generates 8 bits.
%Warning-WIDTHEXPAND: sloppy.sv:15:17: Operator ASSIGNW expects 10 bits on
                      the Assign RHS, but Assign RHS's VARREF 'a'
                      generates 8 bits.
%Warning-ALWCOMBORDER: sloppy.sv:36:15: Always_comb variable driven after
                       use: 'loop'
%Warning-UNUSEDSIGNAL: sloppy.sv:10:15: Bits of signal are not used:
                       'narrow'[2:0]
%Warning-UNUSEDSIGNAL: sloppy.sv:11:15: Signal is not used: 'wide'
%Warning-CASEINCOMPLETE: sloppy.sv:19:5: Case values incompletely covered
                         (example pattern 0x3)
%Warning-BLKSEQ: sloppy.sv:28:7: Blocking assignment '=' in sequential
                 logic process
                 : ... Suggest using delayed assignment '<='
%Error: Exiting due to 10 warning(s)
\end{txtcode}

Two more appear only in slightly different code:

\begin{txtcode}
%Warning-LATCH: sloppy2.sv:4:3: Latch inferred for signal 'sloppy2.lat'
                (not all control paths of combinational always assign
                 a value)
                : ... Suggest use of always_latch for intentional latches
%Warning-UNOPTFLAT: sloppy2.sv:9:15: Signal unoptimizable: Circular
                    combinational logic: 'sloppy2.w'
                    Example path: sloppy2.w
                    Example path: ASSIGNW
                    Example path: sloppy2.w
\end{txtcode}

What each one means to someone coming from \VHDL{}:

\begin{description}
  \item[\code{WIDTHTRUNC} / \code{WIDTHEXPAND}] The assignment silently
        changed the width. In \VHDL{} this is an error at analysis time. Treat
        both as errors and write explicit casts:
        \code{narrow = a[3:0];} or \code{wide = 10'(a);}.
  \item[\code{IMPLICIT}] You used an identifier you never declared and Verilog
        made a one-bit wire out of it. \VHDL{} has no such rule. Put
        \code{`default\_nettype none} at the top of every file and this class
        of typo becomes an error.
  \item[\code{CASEINCOMPLETE}] Not every value of the selector is covered.
        Combined with an \code{always\_comb} that does not assign a default,
        this is a latch. In \VHDL{} the \code{others} choice is compulsory in a
        \code{case}; \SV{} has no such requirement.
  \item[\code{LATCH}] The inferred latch itself. \code{always\_latch} declares
        that you meant it.
  \item[\code{BLKSEQ}] A blocking \code{=} inside \code{always\_ff}. This is
        the \SV{} form of the \VHDL{} variable-versus-signal confusion, and it
        creates simulation-synthesis mismatches.
  \item[\code{UNOPTFLAT}] Verilator found a combinational path from a vector
        back to itself that it cannot split into independent bits. It is
        usually a real design smell and always a performance problem, because
        Verilator must then evaluate that region iteratively.
  \item[\code{UNUSEDSIGNAL}] Something is declared or driven and never read.
        Often harmless, often the fingerprint of a copy-paste error or a
        forgotten connection.
  \item[\code{ALWCOMBORDER}] Inside one \code{always\_comb} a variable is read
        before it is written. The block will still settle, but the intent is
        unclear.
\end{description}

\subsection{Waiving warnings}

\index{lint\_off@\code{lint\_off}}Two mechanisms, and they have different
politics. An inline pragma marks one place in one file:

\begin{svcode}
// verilator lint_off WIDTHTRUNC
assign y = a;
// verilator lint_on WIDTHTRUNC
\end{svcode}

A configuration file, conventionally with the extension \file{.vlt}, keeps the
waivers out of the \RTL{} entirely and is passed on the command line like a
source file:

\begin{txtcode}
// waivers.vlt -- project-wide lint waivers, kept out of the RTL.
`verilator_config
lint_off -rule UNUSEDSIGNAL -file "*sloppy2.sv" -match "*'a'*"
lint_off -rule UNOPTFLAT    -file "*sloppy2.sv" -match "*'sloppy2.w'*"
\end{txtcode}

\begin{shcode}
$ verilator --lint-only -Wall --top-module sloppy2 waivers.vlt sloppy2.sv
%Warning-LATCH: sloppy2.sv:4:3: Latch inferred for signal 'sloppy2.lat'
\end{shcode}

Both targeted warnings are gone and the latch remains, which is what you want.

\begin{tipbox}{Waive in a file, not in the source}
Inline \code{lint\_off} is convenient and therefore dangerous: it spreads, it
is invisible in review, and it is copied along with the code it was written
for. Put waivers in a \file{.vlt} file with a comment saying why, so that the
list of things your project has decided to ignore is one file that a reviewer
can read.
\end{tipbox}

The CORDIC of \refexamplebook passes cleanly, which is the standard to hold
yourself to:

\begin{shcode}
$ verilator --lint-only -Wall --top-module cordic_rot \
      rtl/cordic_pkg.sv rtl/cordic_rot.sv
$ echo $?
0
\end{shcode}

Note that it does so partly by using two narrowly scoped pragmas in
\file{cordic\_pkg.sv} -- \code{lint\_off UNUSEDPARAM} around constants that
only testbenches read, and \code{lint\_off UNUSEDSIGNAL} around the deliberate
truncation in \code{sat\_narrow}. Both are documented in a comment on the line
above. That is the right way to use an inline pragma: narrow, explained, and
immediately turned back on.

\begin{ruleofthumb}
Make \code{verilator -{}-lint-only -Wall} a gate in continuous integration from
the first commit. Retrofitting lint onto a 50\,000-line design is a project;
keeping it clean from the start is free.
\end{ruleofthumb}
% ===========================================================================
\section{Where Verilator Does Not Reach}
\label{sec:verilator:limits}
% ===========================================================================

Four things send you back to \tool{QuestaSim}, and it is worth knowing them
before you build a flow that assumes otherwise.

\textbf{\VHDL{}.} \tool{Verilator} does not read it. A mixed-language
project --- a \VHDL{} block instantiated inside a \SV{} testbench, or the
reverse --- needs a simulator that reads both, and that means a commercial
one. In \tool{QuestaSim} you compile each language with its own compiler,
\code{vcom -2008} and \code{vlog -sv}, into the same library and elaborate the
two together; the port types that survive the boundary are the scalar and
vector logic types, and anything cleverer needs a wrapper.

\textbf{Four-state simulation.} \tool{Verilator} is two-state. A reset bug
whose symptom is an \code{X} propagating through the design will not appear.
\Cref{sec:verilator:twostate} says what to do about it; the honest summary is
that you run the reset tests somewhere else.

\textbf{Covergroups and the constraint solver.} Neither is implemented in
5.020. A testbench that depends on constrained random stimulus and functional
coverage --- which is to say, the one \refverificationbook teaches you to write ---
runs under \tool{Verilator} with both features quietly absent.

\textbf{UVM.} It needs the class library, the DPI back end and language
features \tool{Verilator} does not have. There is no route.

\begin{ruleofthumb}
Lint and regress with \tool{Verilator} because it is fast and free; sign off
with a four-state simulator because your silicon is not two-state.
\end{ruleofthumb}


\section{A Makefile for the Whole Chapter}
\label{sec:verilator:make}

Everything in this chapter is driven from one file, and every target below was
executed to produce the output printed above.

\begin{makecode}
VERILATOR ?= verilator
JOBS      ?= $(shell nproc)

RTL     := rtl/cordic_pkg.sv rtl/cordic_rot.sv
CHK     := cov/cordic_checker.sv
COMMON  := --cc --exe --build -j $(JOBS) --top-module cordic_rot
CFLAGS_O:= -CFLAGS -O2

.PHONY: all first tb trace cov uvm coro svtb lint bench clean

all: first tb cov uvm coro svtb lint

# --- 1. the 40-line first example ------------------------------------
first:
	$(VERILATOR) $(COMMON) -Mdir obj_first $(RTL) tb/main_first.cpp \
	   -o tb_first
	./obj_first/tb_first

# --- 2. the self-checking testbench ----------------------------------
tb:
	$(VERILATOR) $(COMMON) --trace $(CFLAGS_O) -Mdir obj_tb \
	   $(RTL) tb/tb_cordic.cpp -o tb_cordic
	./obj_tb/tb_cordic +n=200

trace:
	$(VERILATOR) $(COMMON) --trace-fst $(CFLAGS_O) -Mdir obj_fst \
	   $(RTL) tb/tb_cordic.cpp -o tb_fst
	./obj_fst/tb_fst +n=200 +trace

# --- 3. coverage and assertions --------------------------------------
cov:
	$(VERILATOR) $(COMMON) --trace --assert \
	   --coverage-line --coverage-toggle --coverage-user \
	   -Mdir obj_cov $(RTL) $(CHK) tb/tb_cordic.cpp -o tb_cov
	./obj_cov/tb_cov +n=300 +bp
	verilator_coverage --annotate annotated logs/coverage.dat

# --- 4. the UVM-shaped C++ environment -------------------------------
uvm:
	$(VERILATOR) $(COMMON) -CFLAGS "-O2 -I../uvmlike" -Mdir obj_uvm \
	   $(RTL) uvmlike/main_uvmlike.cpp -o tb_uvm
	./obj_uvm/tb_uvm +n=200 +stall=25

# --- 5. the coroutine variant ----------------------------------------
coro:
	$(VERILATOR) $(COMMON) -CFLAGS "-std=c++20 -fcoroutines -O2" \
	   -Mdir obj_coro $(RTL) coro/main_coro.cpp -o tb_coro
	./obj_coro/tb_coro +n=200

# --- 6. the SystemVerilog testbench under --binary --timing ----------
svtb:
	$(VERILATOR) --binary --timing -j $(JOBS) $(CFLAGS_O) \
	   --top-module tb_vlt_top -Mdir obj_svtb \
	   rtl/cordic_pkg.sv svtb/cordic_vif.sv rtl/cordic_rot.sv \
	   svtb/tb_vlt_pkg.sv svtb/tb_vlt_top.sv -o tb_svtb
	./obj_svtb/tb_svtb +n=200

# --- 7. lint ----------------------------------------------------------
lint:
	$(VERILATOR) --lint-only -Wall --top-module cordic_rot $(RTL)
	@echo "lint clean"

bench:
	./obj_tb/tb_cordic +n=200000

clean:
	rm -rf obj_* annotated logs *.vcd *.fst
\end{makecode}

\begin{tipbox}{Give every variant its own \code{-Mdir}}
Coverage, tracing and \code{-{}-timing} all change the generated code, so a
single \file{obj\_dir} would be rebuilt from scratch every time you switch
targets. Separate directories cost disk and save minutes.
\end{tipbox}

\section{Checklist}
\label{sec:verilator:checklist}

\begin{itemize}
  \item Verilator is a \emph{compiler}. The \DUT{} becomes a C++ class, the
        testbench is a \code{main()}, and there is no interactive prompt, no
        event queue for most signals, and no \code{\$time} unless you keep one.
  \item \textbf{Never trust a Verilator run to catch a missing reset.} It is
        two-state and X-optimistic. Run reset checks in a four-state simulator
        (\tool{QuestaSim}, \tool{Icarus}), or run twice with
        \code{+verilator+rand+reset+2} and different
        \code{+verilator+seed+} values and require identical output.
  \item \code{eval()} means ``settle'', not ``advance one cycle''. A clock edge
        is you writing \code{clk\_i} and calling \code{eval()} again. Sample
        handshake signals \emph{between} the settling \code{eval()} and the
        edge.
  \item \code{ctx.timeInc()} counts \emph{timeprecision} units. Under
        \code{`timescale 1ns/1ps} a 5\,ns half period is \code{timeInc(5000)}.
  \item Every port is unsigned in C++. Cast a signed output to \code{int16\_t}
        or \code{int32\_t} yourself; the \SV{} sign information does not cross
        the boundary. Widths map \code{CData}/\code{SData}/\code{IData}/
        \code{QData}/\code{WData} at 8/16/32/64/65+ bits.
  \item \code{-{}-trace} or \code{-{}-trace-fst} must be given at verilation
        time as well as switched on at run time; otherwise you get a link
        error, not a message about tracing. \code{traceEverOn(true)} comes
        first.
  \item In a \code{-{}-cc} build nothing writes the coverage database for you:
        call \code{ctx.coveragep()->write("logs/coverage.dat")} before exit.
  \item In 5.020: classes, mailboxes, semaphores, \code{fork}/\code{join*},
        virtual interfaces, \code{bind} and DPI-C all work. Covergroups,
        constraint solving, clocking blocks through a virtual interface,
        \code{\#\#} sequence delays, SDF and UVM do not.
  \item \code{randomize()} returns 1 and ignores your constraints, after one
        \code{CONSTRAINTIGN} warning at build time. Guard the class with
        \code{`ifdef VERILATOR} and supply a hand-written randomiser.
  \item Never start a comment with the word ``Verilator'': it is parsed as a
        metacomment and becomes a syntax error.
  \item When you need UVM's structure and cannot have UVM, build the four
        things it is actually for: a named component tree with phases, an
        analysis port, a factory with type overrides, and a configuration
        store. Two hundred lines of C++ is enough.
  \item Blocking tasks do not exist in C++. Use \code{tick()} state machines
        first; use C++20 coroutines only when the state machine has become
        unreadable, and expect to pay about 40\% in speed and much more in
        lifetime bugs.
  \item Run \code{verilator -{}-lint-only -Wall} on every commit. It is the
        cultural replacement for the \VHDL{} type checker
        (\refdesignbook). Put \code{`default\_nettype none} at the top of
        every file, treat \code{WIDTHTRUNC}, \code{WIDTHEXPAND},
        \code{CASEINCOMPLETE}, \code{LATCH} and \code{BLKSEQ} as errors, and
        keep waivers in a \file{.vlt} file rather than inline.
  \item Verilator wins on long runs and loses on short ones: verilation of the
        CORDIC took 0.08\,s and the C++ build 5.5\,s, so a two-second
        regression is not where the tool pays for itself.
  \item \textbf{Four things send you back to \tool{QuestaSim}}: \VHDL{} and
        mixed-language elaboration, four-state semantics, covergroups and the
        constraint solver, and UVM. \Cref{sec:verilator:limits} says what to do
        about each.
\end{itemize}
